% Converted on 2026-10-08 from archive/2026-10-08-development-beamer/appendices/b-decisions-references.tex (Beamer, review copy 43).
% Each frame became a subsection whose body, unchanged, sits in a group as in Beamer.

\section{결정·관리·참고자료}

\hypertarget{appendix-b}{}

\subsection{결정 기록: 합의와 제안을 구분한다}\label{decisions}
\begingroup% frame-local settings, as in Beamer
  \framesubtitle{기록일 2026-09-23 · 갱신 2026-09-29 · 대화의 합의·결정은 세부 API 구현 완료와 다르다.}
  \small\renewcommand{\arraystretch}{1.12}
  \begin{tabularx}{\textwidth}{@{}p{0.7cm}YP{2.5cm}@{}}
    \toprule
    ID & 내용 & 상태 \\
    \midrule
    D01 & 같은 물리 모델을 여러 계산 방법으로 전달하는 workflow & \textbf{방향 합의} \\
    D02 & 모델 간 공통 필드와 동일한 의미의 전달 규약 & \textbf{방향 합의} \\
    D03 & 모델 작업 공간은 작은 구성으로 시작 & \textbf{방향 합의} \\
    D04 & 원시·중간 계산은 모델 작업 공간, 정돈된 결과는 data-space & \textbf{방향 합의} \\
    D05 & 분율 좌표·정수 셀 이동·단위 모드·결합 수 규칙 & \textbf{규약 검토} \\
    D06 & 기존 변위와 새 규약의 대응, 실제 자료형 API & \textbf{D13·D16으로 결정} \\
    D07 & 1차 범위: NBCP·두 하이젠버그 벤치마크, LSWT 물리량(위상량은 마지막 단계) & \textbf{사용자 결정} \\
    D08 & \texttt{terms[]} 형식: 형식은 열고 1차 검증은 이선형·Zeeman으로 닫음 & \textbf{사용자 결정} \\
\bottomrule
  \end{tabularx}
  \par\smallskip
  \supportingtext{후속 변경 시 각 항목의 결정 근거와 검증을 갱신한다.
  구현·수치 검증·사용자 확인은 각각 별도 상태로 기록한다.}
  \slidesource{대화 합의와 검토안의 범위}{contract}
\endgroup


\subsection{결정 기록 (계속)}\label{decisions-continued}
\begingroup% frame-local settings, as in Beamer
  \framesubtitle{2026-09-29 사용자 결정.}
  \small\renewcommand{\arraystretch}{1.12}
  \begin{tabularx}{\textwidth}{@{}p{0.7cm}YP{2.5cm}@{}}
    \toprule
    ID & 내용 & 상태 \\
    \midrule
    D09 & 모델은 g-텐서, 장 $B$는 외부 변수; Zeeman 부호 $-\mu_B$ & \textbf{사용자 결정} \\
    D10 & 최소 ED 검증 도구를 \texttt{methods/ed/}에 두고 ED 솔버로 확장 & \textbf{사용자 결정} \\
    D11 & 패키지 \texttt{spintoolkit}, 별칭 \texttt{stk}; 키타에프는 후속 벤치마크 & \textbf{사용자 결정} \\
    D12 & 선택적 Zeeman 항, $g_{\alpha\beta}$ 인덱스, relative $\mu_B=1$(D20으로 대체), 상태 검증·기준 상태 진단 & \textbf{사용자 결정} \\
    D13 & 표준 Fourier 부호 $e^{+\mathrm i\mathbf k\cdot\mathbf r}$와 전체 위치 게이지; NBCP 변위는 $r_t=r_s-d$ & \textbf{사용자 결정} \\
    D14 & 상태·계산 요청·결과를 공통·방법별 부분으로 분리; \texttt{SpinState}는 정수 초격자 $M$ & \textbf{사용자 결정} \\
    D15 & 표준 벤치마크 모델은 \texttt{spintoolkit/models/}, 연구 모델은 \texttt{code-space/model/<name>/} & \textbf{사용자 결정} \\
    D16 & 1단계 API: \texttt{Term} 생성 함수, (사이트, 셀) 방향 dict, 위반 일괄 보고, 최상위 노출; 직렬화는 4단계 & \textbf{사용자 결정} \\
    \bottomrule
  \end{tabularx}
  \par\smallskip
  \supportingtext{구현·수치 검증·사용자 확인은 각각 별도 상태로 기록한다.}
  \slidesource{2026-09-29 결정}{contract}
\endgroup


\subsection{결정 기록 (D17--D24)}\label{decisions-selection}
\begingroup% frame-local settings, as in Beamer
  \framesubtitle{2026-09-29--30 사용자 결정. 근거와 수치 자료는 부록 A와 이 부록의 뒤쪽에 있다.}
  \footnotesize\renewcommand{\arraystretch}{1.08}
  \begin{tabularx}{\textwidth}{@{}p{0.7cm}YP{2.5cm}@{}}
    \toprule
    ID & 내용 & 상태 \\
    \midrule
    D17 & 영점 에너지 상태 선택은 고전 manifold 위에서: $E_{\mathrm{cl}}$ Hessian 영공간으로 정한 축 $\hat n$의 궤도 $R_{\hat n}(\phi)$에서 $E_{\mathrm{qm}}$ 최소화(2b단계) & \textbf{사용자 결정} \\
    D18 & D17 구현 계획: \texttt{quantum} 삭제(과거의 실수), MAGSWT 유지·재현 검증, 2b단계, 삼각격자 $h=0$·$h>h_{\mathrm{sat}}$ & \textbf{사용자 결정} \\
    D19 & D17 판정: 물리 근거 모드(기본)와 고정 상수 모드; 기본값은 \texttt{definitions/defaults.py}; 판정 값 기록 & \textbf{사용자 결정} \\
    D20 & 수치 계산은 무차원($E_0$ 기준: 장 $\mu_B\mathbf B/E_0$, 온도 $k_BT/E_0$); 물리 단위 T·K·meV는 문서 규약, 변환은 사용자 & \textbf{사용자 결정} \\
    D21 & 2단계: NBCP 문헌·원고 파라미터 세트(g는 세트별), 변환은 \texttt{system/conversion.py}, 장이 같을 때만 역변환, 새 LSWT 진입점은 4단계 & \textbf{사용자 결정} \\
    D22 & 2b단계: \texttt{select\_on\_manifold}; 경쟁 영역은 두 최소를 기록하고 자동 선택 안 함; $E_{\mathrm{qm}}$은 호출 가능 객체(기존 \texttt{EnergyFunction} 경유); 새 DE 없음; \texttt{quantum} 삭제 & \textbf{사용자 결정} \\
    D23 & 3단계: 일반 해밀토니안 ED(U(1)·운동량 섹터는 선택), 2-마그논·flux는 후속, 문헌 비교는 4단계; 토러스의 겹친 결합 합산, 접히는 결합은 모든 방법에서 거부; ED는 전체 에너지 저장 & \textbf{사용자 결정} \\
    D24 & 4단계: \texttt{solve\_lswt} 진입점(기존 \texttt{LSWTSolver} 유지), 무차원, 정규화 기본 없음(영모드는 오류), 대각화 정보 보관·재사용, 4a·4b·4c & \textbf{사용자 결정} \\
    \bottomrule
  \end{tabularx}
  \par\smallskip
  \supportingtext{구현·수치 검증·사용자 확인은 각각 별도 상태로 기록한다.}
  \slidesource{2026-09-29 결정}{contract}
\endgroup


\subsection{결정 기록 (D25--D28)}\label{decisions-stage4}
\begingroup% frame-local settings, as in Beamer
  \framesubtitle{2026-09-30 사용자 결정. 근거와 수치 자료는 부록 A에 있다.}
  \small\renewcommand{\arraystretch}{1.12}
  \begin{tabularx}{\textwidth}{@{}p{0.7cm}YP{2.5cm}@{}}
    \toprule
    ID & 내용 & 상태 \\
    \midrule
    D25 & 4b: 영모드는 $\lambda_{\min}(\mathsf H)/\text{scale}$로 영($\le10^{-12}$)·후보($\le10^{-6}$)·갭; 후보가 있으면 사용자가 \texttt{gapless}를 정할 때까지 유한 온도 계산을 멈춤; gapless이면 $t>0$의 보손 수·모멘트·자화는 NaN, $F,U,S,C$는 계산; 탐색은 고대칭점·격자 최소점·연속 최소화·선 검사, $k=0$ 기원 표시; 무차원 $t$ & \textbf{사용자 결정} \\
    D26 & 4c: 한 마그논 가로 성분과 정확한 Bragg(2-마그논 연속체는 후속); 모드별 무게가 주 출력, 넓힌 스펙트럼은 보조; 새 구현, 기존 \texttt{correlations.py}는 비교 대상만; 결합 상관과 에너지 교차 확인 & \textbf{사용자 결정} \\
    D27 & 기존 코드 정리: \texttt{correlations.py} 삭제, 동시간 실공간 상관·사다리 성분은 새 모듈로; MAGSWT 격자 탐색 삭제, 궤도 위 에너지 지형은 \texttt{orbit\_energy\_landscape}; MAGSWT 정규화는 유지 & \textbf{사용자 결정} \\
    D28 & 고전·양자 비교: 완화 경로 위 1-loop 유효 퍼텐셜 $\Gamma(\phi)=E_{\mathrm{cl}}+E_{\mathrm{zp}}$의 최소(부드러운 방향의 1-loop 정상성); 곡률 비 대역을 대체; 평평하면 D17 & \textbf{사용자 결정}, 구현·제공자 결정 \\
    \bottomrule
  \end{tabularx}
  \slidesource{2026-09-30 결정}{contract}
\endgroup


\subsection{결정 기록 (D29)}\label{decisions-stage5}
\begingroup% frame-local settings, as in Beamer
  \framesubtitle{2026-09-30 사용자 결정(추천안 채택). 근거와 수치 자료는 부록 A에 있다.}
  \small\renewcommand{\arraystretch}{1.12}
  \begin{tabularx}{\textwidth}{@{}p{0.7cm}YP{2.5cm}@{}}
    \toprule
    ID & 내용 & 상태 \\
    \midrule
    D29 & 5단계 위상량: thermal Hall은 층당 $\kappa_{xy}^{2D}/T$를 $k_B^2/\hbar$ 단위의 무차원 값으로(5b); 곡률은 D13 전체 위치 게이지에서 계산·보고, Chern 수·$\kappa$의 게이지 불변은 시험; 부호 $\Omega=\partial_{k_x}A_y-\partial_{k_y}A_x$, $A=i\langle u|\eta\nabla u\rangle$, $C=\frac{1}{2\pi}\int\Omega$, $\kappa=-\frac{k_B^2T}{\hbar A}\sum c_2\Omega$, 물리적 Bloch 상태(ED)로 고정; 이상항 기준 모델로 키타에프 [111] 편극상; 5a·5b·5c로 나눠 검증. 5b 세부: $\kappa$는 쌍 합(퇴화·교차 밴드에서도 정의), 영모드는 D25. 5c: 두 Hall 경로의 계산 핵심 통합, 기존 SI API 유지. 5d: 적응형 k 적분($\kappa$만, 선택 인자) & \textbf{사용자 결정}, 5a--5d 구현 \\
    \bottomrule
  \end{tabularx}
  \slidesource{2026-09-30 결정}{contract}
\endgroup


\subsection{결정 기록 (D30): 설계 미결 사항}\label{decisions-d30}
\begingroup% frame-local settings, as in Beamer
  \framesubtitle{2026-09-30 사용자 결정(추천안 채택). 구현은 LSWT 작업 재개 뒤에 한다.}
  \footnotesize\renewcommand{\arraystretch}{1.1}
  \begin{tabularx}{\textwidth}{@{}p{2.4cm}Y@{}}
    \toprule
    항목 & 결정 \\
    \midrule
    자기 셀 & LT는 진단 도구로 구현($J(\mathbf q)$ 최소 $\mathbf q^*$, 강한 제약 만족 여부, 후보 셀 제안; 선택 단계는 그대로). 대각 셀 전용 \texttt{CommensurateStructure}는 삭제(참조·호환 매핑·예시·테스트 정리, 옛 저장 객체 복원 여부 먼저 확인) \\
    기존 경로 & \texttt{LSWTSolver}($\to$ \texttt{solve\_lswt}), SI Hall API($\to$ \texttt{thermal\_hall})에 사용 중단 경고, 삭제는 공개 배포 정리 때. 전역 고전 탐색을 새 자료형으로 옮긴 뒤 \texttt{EnergyFunction}/\texttt{SpinOptimizer}도 같은 방식으로. \texttt{SpinSystem}과 기존 해밀토니안 구성기는 내부 도구로 유지 \\
    1차 범위 제외 & D17 확장(사이트별 회전축·여러 차원 영공간), 후속 항 종류(같은 사이트 이차항, 3·4-스핀 항, $S\ge1$), 키타에프 flux 섹터 대조와 ED 2-마그논 섹터 --- 필요한 벤치마크나 모델이 생길 때 다시 연다 \\
    TN 필드 규약 & 솔버 경계 검증 실험(인수인계 260607) 결과를 보고 정한다 \\
    적응형 적분 & 오차 추정 개선은 보류(계산량이 문제가 될 때) \\
    \bottomrule
  \end{tabularx}
  \slidesource{2026-09-30 결정}{contract}
\endgroup


\subsection{결정 기록 (D31): 위상량 수치 기준 확정}\label{decisions-d31}
\begingroup% frame-local settings, as in Beamer
  \framesubtitle{2026-09-30 사용자 결정. 두 계산은 반대 방향으로 실패하므로, 일치할 때만 받으면 틀린 정수 대신 NaN이 나온다.}
  \footnotesize\renewcommand{\arraystretch}{1.1}
  \begin{tabularx}{\textwidth}{@{}P{3.2cm}P{1.6cm}Y@{}}
    \toprule
    기준 & 값 & 의미와 근거 \\
    \midrule
    밴드 간격 & $10^{-8}E_0$ & 이하이면 그 격자점의 곡률 판정 불가(Kubo 분모가 0에 가까움). 기존 경로의 승인값을 무차원으로 \\
    FHS 링크 겹침 & $10^{-3}$ & 이하이면 격자점 사이의 밴드 교차. 정상 경우 최소 겹침 0.55--0.95, 대칭 보호 교차 $10^{-15}$ \\
    Kubo--FHS 일치 & 0.1 & FHS: 갭이 있으면 거친 격자에서도 정확, 갭 닫힘은 모름(임의의 정수). Kubo: 작은 갭에서 느리게 수렴, 갭 닫힘·격자 부족이 정수 이탈로 드러남. 거부되면 격자를 키운다 \\
    FHS admissibility & $10^{-6}$ rad & plaquette 위상이 $\pm\pi$에서 이 이내이면 판정 불가. 밴드 접촉을 둘러싼 plaquette의 Berry 위상 $\pi$는 가지($\pm\pi$)가 반올림으로 정해짐. 최대 위상: $D=0$에서 $\pi$, $D=10^{-3}$에서 $\pi-0.06$, 갭 있는 경우 1.01 이하 (2026-10-01 보완) \\
    적응형 적분 & 1e-3 / 1e-7 / 2e5 / 12 & 상대·절대 허용오차, 계산점 예산, 최대 깊이 \\
    \bottomrule
  \end{tabularx}
  \par\smallskip
  \supportingtext{Haldane 예: $D=10^{-2}$--$10^{-4}$에서 FHS는 $6\times6$부터 정답 $(+1,-1)$, Kubo는 0.002--0.99로 수렴 중이라 대부분 거부; $D=0$(갭 없음)에서 FHS는 플랫폼·격자에 따라 $-1$, 0, $+1$이었고, 0이면 Kubo 0과 일치해 틀린 $(0,0)$이 채택됐다($24\times24$) $\to$ 2026-10-01 admissibility 기준으로 판정 불가.}
  \slidesource{2026-09-30 결정}{contract}
\endgroup


\subsection{결정 기록 (D32--D33): 6단계와 4d 인터페이스}\label{decisions-d32}
\begingroup% frame-local settings, as in Beamer
  \framesubtitle{2026-09-30 사용자 결정(추천안 채택). D30의 구현 방법.}
  \small\renewcommand{\arraystretch}{1.12}
  \begin{tabularx}{\textwidth}{@{}P{3.0cm}Y@{}}
    \toprule
    항목 & 결정 \\
    \midrule
    LT 진단 & \texttt{luttinger\_tisza(model, mesh=(N, N)) -> LTReport}(\texttt{methods/luttinger\_tisza.py}): $\mathbf q^*$, 분수와 후보 초격자, 고유공간 차원, 단일 q 강한 제약과 그 상태, 퇴화 표시; 영장; 선택하지 않음 \\
    전역 고전 탐색 & \texttt{classical\_search(model, supercell, conditions, seed=42, ...)}(\texttt{methods/classical.py}): $(\theta,\phi)$ 차분 진화(기존 설정) 뒤 \texttt{refine\_classical}; 각도 고정 없음 \\
    진행 & 6a(정리) $\to$ 6b $\to$ 6c; 6b·6c는 검증 뒤 사용자 확인 \\
    \bottomrule
  \end{tabularx}
  \slidesource{2026-09-30 결정}{contract}
  \par\medskip
  \begin{tabularx}{\textwidth}{@{}P{3.0cm}Y@{}}
    \toprule
    D33 & 결정 \\
    \midrule
    밴드 구조(4d) & \texttt{high\_symmetry\_points(lattice)}, \texttt{band\_structure(result, path, points, lattice="primitive")}: 결정 격자 경로 기본, 자기 셀 밴드는 접힌 채, 영모드 표시, 불안정은 NaN; 그림은 7단계(시각화)에서 설계·검토 \\
    \bottomrule
  \end{tabularx}
\endgroup


\subsection{결정 기록 (D34): 단일 Q 나선의 회전틀 LSWT}\label{decisions-d34}
\begingroup% frame-local settings, as in Beamer
  \framesubtitle{2026-10-01 Claude 결정(사용자 위임), 2026-10-13 사용자 승인(일괄 검토 1번). Toth--Lake (2015).}
  \small\renewcommand{\arraystretch}{1.12}
  \begin{tabularx}{\textwidth}{@{}P{3.0cm}Y@{}}
    \toprule
    항목 & 결정 \\
    \midrule
    상태 & \texttt{IncommensurateStructure(model\_ref, wave\_vector, rotation\_axis, directions)}: $\mathbf n_a(\mathbf R)=R_{\mathbf n}(2\pi\,\mathbf q\cdot(n_1,n_2))\,\mathbf d_a$, 위상은 셀 번호만, 원뿔 허용, 정합이면 \texttt{to\_spin\_state} \\
    LSWT & \texttt{solve\_spiral\_lswt}: $J'=J R_{\mathbf n}(2\pi\,\mathbf q\cdot\Delta\mathbf n)$인 원시 셀 모델을 \texttt{solve\_lswt}로 풂; 밴드는 회전틀 운동량의 $\omega(\mathbf k)$ \\
    적용 조건 & 모든 $J$가 축 회전과 교환(축 XXZ, 축 방향 DM), $g^T\mathbf b\parallel\mathbf n$; 아니면 근사값 대신 \texttt{SpiralSymmetryError} \\
    구조인자 & \texttt{spiral\_structure\_factor}: $S=\sum_m R_m S'(\mathbf k-m\mathbf Q)R_m^\dagger$; $\mathbf Q$ 또는 $2\mathbf Q\in G$이면 초격자 안내 \\
    안전장치 & 회전틀 결과는 실험실 상관·위상량이 거부(나선 위상량 미검증); 피치 $\partial E/\partial\mathbf q$ 진단 \\
    LT 연결 & \texttt{IncommensurateStructure.from\_lt}, \texttt{LTWaveVector.amplitude}, \texttt{refine\_spiral} \\
    \bottomrule
  \end{tabularx}
  \par\smallskip
  \supportingtext{검증: 삼각 $120^\circ$·DM·원뿔 나선이 같은 운동량의 초격자 LSWT와 에너지·밴드·$S^{ab}(\mathbf q,\omega)$에서 반올림 수준 일치, 비정합 J1--J2 사슬의 피치와 $\omega(\mathbf k)$가 해석식과 $10^{-14}$ 수준 일치 (\texttt{verification/spiral-rotating-frame-2026-10-01.json}).}
  \slidesource{2026-10-01 나선 메모}{contract}
\endgroup


\subsection{결정 기록 (D35): 결정 대칭과 대칭 허용 결합}\label{decisions-d35-symmetry}
\begingroup% frame-local settings, as in Beamer
  \framesubtitle{2026-10-01 Claude 결정(사용자 위임: 범용 솔버 방향), 2026-10-13 사용자 승인(일괄 검토 2번). 기존 API 변경 없음.}
  \small\renewcommand{\arraystretch}{1.12}
  \begin{tabularx}{\textwidth}{@{}P{3.0cm}Y@{}}
    \toprule
    항목 & 결정 \\
    \midrule
    위치·입력 & \texttt{system/symmetry.py}; 결정은 \texttt{LayerCrystal(lattice, positions, species, heights)}로 리간드까지 받는다. 자기 이온만의 대칭(\texttt{LayerCrystal.from\_model})은 상한이다 \\
    연산 & $\mathbf r\to R\mathbf r+\mathbf t$, $R=\mathrm{diag}(R_2,\pm1)$; 스핀·장은 축 벡터 $R_s=\det(R)R$; $J\to R_sJR_s^{T}$, 결합이 뒤집히면 전치; $g\to R_sgR_s^{T}$; 시간 반전은 지원 항을 모두 보존하므로 군에 넣지 않음 \\
    기능 & \texttt{find\_symmetry}, \texttt{close\_group}; \texttt{CrystalSymmetry.allowed\_exchange}, \texttt{allowed\_g\_tensor}(허용 형태의 기저), \texttt{bilinear\_terms}, \texttt{zeeman\_terms}(궤도 생성), \texttt{check\_model} \\
    물리 판정 & 대칭을 깨는 계수는 금지 성분을 보여 주며 거부(임의로 대칭화하지 않음) \\
    검증 & 평평한 삼각격자 XYZ 3개, CdI$_2$형 D$_{3d}$ 삼각격자 $J,\Delta,J_{\pm\pm},J_{z\pm}$ 4개·DM 없음·축 대칭 $g$, 평평한 카고메 $D_z$만·균일 순환; 생성 모델의 LSWT $\omega(C_3\mathbf k)=\omega(\mathbf k)$ \\
    \bottomrule
  \end{tabularx}
  \slidesource{2026-10-01 결정}{contract}
\endgroup


\subsection{결정 기록 (D36--D37): 직접 $H(\mathbf k)$와 단일 이온 항}\label{decisions-d36-onsite}
\begingroup% frame-local settings, as in Beamer
  \framesubtitle{2026-10-01 Claude 결정(사용자 위임; 사용자가 2번 집중을 지정); D37은 2026-10-13 사용자 승인(일괄 검토 3번), D36은 2026-10-15 사용자 승인(일괄 검토 목록 밖 추가 항목).}
  \small\renewcommand{\arraystretch}{1.12}
  \begin{tabularx}{\textwidth}{@{}P{3.0cm}Y@{}}
    \toprule
    항목 & 결정 \\
    \midrule
    D36 직접 $H(\mathbf k)$ & \texttt{methods/lswt/quadratic.py}의 \texttt{QuadraticBoseHamiltonian}: \texttt{SpinModel}에서 $H(\mathbf k)$와 해석적 $\partial_{\mathbf k}H$를 바로 만든다. \texttt{solve\_lswt}가 사용; 규약(Nambu 순서, 국소 좌표 게이지, D13)은 그대로이며 기존 구성기와 $10^{-16}$ 수준으로 일치 \\
    D37 항 & \texttt{Term.onsite(site, A)}: $\mathbf S_i^{T}A\mathbf S_i$, $A$는 실수 대칭(반대칭 부분은 비에르미트라 거부), 사이트당 하나 \\
    D37 규칙 & 고전·LSWT는 스핀 결맞음 상태 값 $S(S-\tfrac12)\mathbf n^{T}A\mathbf n+\tfrac S2\operatorname{tr}A$, 즉 계수 $(1-\tfrac1{2S})A$를 쓰고, 횡방향 정규 순서 상수는 더하지 않는다. ED와 LT는 각각 정확한 연산자와 같은 결맞음 값을 쓴다 \\
    근거 & $S=\tfrac12$에서 상수 $\operatorname{tr}A/4$; $D(S^z)^2$ 단일 이온 갭 $(2S-1)|D|$ 정확; 선형항이 정확한 연산자와 같음; 단일 스핀 $S\le20$ 정확 스펙트럼과 비교 시 large-S는 갭 오차 $\approx0.57$, 정규 순서 상수를 더하면 에너지 오차가 $S$에 비례 \\
    검증 & 축 대칭 모델의 1-마그논 LSWT 밴드 = 토러스 ED($10^{-11}$); 결합 상관 에너지 = LSWT 바닥 에너지; LT 하한 \\
    \bottomrule
  \end{tabularx}
  \slidesource{2026-10-01 결정}{contract}
\endgroup


\subsection{결정 기록 (D38): 고전 스핀 동역학}\label{decisions-d38-dynamics}
\begingroup% frame-local settings, as in Beamer
  \framesubtitle{2026-10-01 Claude 결정(사용자 위임; 사용자가 5번 집중을 지정), 2026-10-14 사용자 잠정 승인(일괄 검토 4번; 동역학은 추후 다시 검토).}
  \small\renewcommand{\arraystretch}{1.12}
  \begin{tabularx}{\textwidth}{@{}P{3.0cm}Y@{}}
    \toprule
    항목 & 결정 \\
    \midrule
    계 & \texttt{methods/dynamics.py}의 \texttt{ClassicalTorus}: 유한 토러스(D23), 스핀 길이 $S_i$, 고전 에너지는 D37 결맞음 값. 시간 단위 $\hbar/E_0$ \\
    운동 방정식 & $\dot{\mathbf s}_i=\mathbf s_i\times\mathbf h_i$, $\mathbf h_i=-\partial E/\partial\mathbf s_i$. 암시적 중점법은 길이와 에너지를 정확히 보존하고 작은 진동 진동수를 $(2/\Delta t)\arctan(\omega\Delta t/2)$로 옮긴다 \\
    랑주뱅 & $\dot{\mathbf s}=\mathbf s\times(\mathbf h+\boldsymbol\xi)-(\alpha/S)\,\mathbf s\times(\mathbf s\times\mathbf h)$, $\langle\xi\xi\rangle=(2\alpha T/S)\delta$ (Stratonovich, 확률 Heun 후 투영). 감쇠 이동도 $\alpha S$와 구면 확산 $\alpha TS$의 비가 $T$라 정상 분포가 $e^{-E/T}$ \\
    $S(\mathbf q,\omega)$ & 완전 위치 위상(D13), Hann 창 주기도, 합 규칙 보존. 양자 상세 균형 인자와 형태 인자는 적용하지 않음 \\
    검증 & 단일 스핀 Langevin 함수; 단일 이온 분포(구적법, large-S 값과 구별); 작은 진동 = LSWT 밴드; 저온 $S(\mathbf q,\omega)$ 정점 = LSWT 분산 \\
    \bottomrule
  \end{tabularx}
  \slidesource{2026-10-01 결정}{contract}
\endgroup


\subsection{결정 기록 (D39): 유한 온도 각도 자유 에너지}\label{decisions-d39-free-energy}
\begingroup% frame-local settings, as in Beamer
  \framesubtitle{2026-10-01 Claude 결정(사용자 위임; 사용자가 NBCP 우선을 지정), 2026-10-14 사용자 잠정 승인(일괄 검토 5번; 실제 계산에 쓸 때 다시 검토).}
  \small\renewcommand{\arraystretch}{1.12}
  \begin{tabularx}{\textwidth}{@{}P{3.0cm}Y@{}}
    \toprule
    항목 & 결정 \\
    \midrule
    식 & $f(\phi,T)=E_{\rm cl}+E_{\rm zp}+\frac{T}{N}\sum_{\mathbf k n}\ln(1-e^{-\varepsilon_{\mathbf k n}/T})$, \texttt{LSWTHarmonicFreeEnergy}; 제한 상태와 영점 처리는 D28 그대로. 1/S 조화 차수 \\
    부드러운 모드 & $|\mathbf k|<\Lambda$의 가장 낮은 모드는 위상 이론(clock RG)에 남기고 열항에서 뺀다. $\Lambda=0$은 전체 조화 자유 에너지. $\Lambda$와 뺀 모드 수를 기록 \\
    정의되지 않음 & 불안정 모드 쌍이나 0 에너지 모드가 남으면 \texttt{nan}과 횟수 기록 \\
    검증 & 편극 정사각 반강자성체 닫힌 식; $-\partial F/\partial T$ = Bose 엔트로피; 차단 모드 차이; NBCP Y·V 메쉬 수렴 \\
    \bottomrule
  \end{tabularx}
  \slidesource{2026-10-01 결정}{contract}
\endgroup


\subsection{결정 기록 (D40): 고전 몬테카를로와 열 관측량}\label{decisions-d40-monte-carlo}
\begingroup% frame-local settings, as in Beamer
  \framesubtitle{2026-10-01 Claude 결정(사용자 위임; 사용자가 NBCP 우선을 지정), 2026-10-14 사용자 잠정 승인(일괄 검토 6번; 실제 계산에 쓸 때 다시 검토).}
  \small\renewcommand{\arraystretch}{1.12}
  \begin{tabularx}{\textwidth}{@{}P{3.0cm}Y@{}}
    \toprule
    항목 & 결정 \\
    \midrule
    갱신 & 결합 그래프 색칠로 같은 색을 한꺼번에; Metropolis 제안 $S\,\mathrm{normalize}(\mathbf s/S+\sigma\mathbf g)$(대칭), onsite 포함 정확한 $\Delta E$; 과완화(국소장 반사)는 onsite 없는 사이트만 \\
    관측량 & 에너지, 비열 $N\,\mathrm{var}(e)/T^2$, $\mathbf m_{\mathbf q}$(D13), 상대 비틀림 강성 $\Upsilon=(\langle E''\rangle-\mathrm{var}(E')/T)/N$, 잭나이프 오차. 축에 대해 U(1)이 아닌 교환은 거부하고, 요청 시 U(1) 대칭 부분 $\bar J$의 강성만 준다 \\
    검증 & 비틀림 미분 = 유한 차분; $\Upsilon(0)=|J|S^2$; 과완화 에너지 보존; 단일 스핀·단일 이온 분포; 고온 급수; 몬테카를로 = Langevin 에너지 \\
    \bottomrule
  \end{tabularx}
  \slidesource{2026-10-01 결정}{contract}
\endgroup


\subsection{결정 기록 (D41): 비편극 중성자 세기}\label{decisions-d41-neutron}
\begingroup% frame-local settings, as in Beamer
  \framesubtitle{2026-10-01 Claude 결정(사용자 위임; NBCP 재현과 배포의 공통 1순위), 2026-10-14 사용자 승인(일괄 검토 7번, 추천대로).}
  \small\renewcommand{\arraystretch}{1.12}
  \begin{tabularx}{\textwidth}{@{}P{3.0cm}Y@{}}
    \toprule
    항목 & 결정 \\
    \midrule
    정의 & 사이트당 $I(\mathbf Q,\omega)=\sum_{ab}(\delta_{ab}-\hat Q_a\hat Q_b)\,S_M^{ab}(\mathbf Q,\omega)$, $S_M$은 $\mathbf S_i\to\tfrac12 g_i F_i(|\mathbf Q|)\mathbf S_i$로 바꾼 D26 구조인자(탄성 포함, 나선 D34 포함). 계수 $(\gamma r_0)^2 k_f/k_i\,e^{-2W}$는 제외 \\
    2D 층 & $S_M$은 $\mathbf Q_\parallel$에만 의존; $Q_z$는 편극 인자와 형상인자의 $|\mathbf Q|$에만 들어간다 \\
    형상인자 & 쌍극자 근사 $\langle j_0\rangle+c\langle j_2\rangle$, Brown 표 중 Sunny.jl·periodictable이 일치하는 3d 이온; $c=0$ 기본 \\
    분해능·평균 & 에너지 의존 FWHM(Gaussian/Lorentzian); 분말 = 구면 Fibonacci 평균; 도메인 $I_R(\mathbf Q)=I(R^{-1}\mathbf Q)$, $R$의 해밀토니안 대칭성은 사용자 책임 \\
    검증 & $g=2,F=1$ 환원; 강자성 $W(1+\hat Q_z^2)$, 비등방 $g$; 나선 = 초격자; 분말 $4/3$; C4 도메인; NBCP 3.5 T 닫힌 해와 $10^{-15}$ 일치 \\
    \bottomrule
  \end{tabularx}
  \slidesource{2026-10-01 결정}{contract}
\endgroup


\subsection{결정 기록 (D42): 스커미온 수와 후보 상태 비교}\label{decisions-d42-texture}
\begingroup% frame-local settings, as in Beamer
  \framesubtitle{2026-10-01 Claude 결정(사용자 위임; 두 목표 공통 2순위), 2026-10-15 사용자 승인(일괄 검토 8번; 후보 비교의 안정성 검사는 D53에서 연속 k 검사 D51을 기본으로 씀).}
  \small\renewcommand{\arraystretch}{1.12}
  \begin{tabularx}{\textwidth}{@{}P{3.0cm}Y@{}}
    \toprule
    항목 & 결정 \\
    \midrule
    스커미온 수 & 주기 삼각분할의 반시계 삼각형마다 Berg--L\"uscher $\tan(\Omega/2)=\mathbf n_i\cdot(\mathbf n_j\times\mathbf n_k)/(1+\sum\mathbf n\cdot\mathbf n)$, $Q=\sum\Omega/4\pi$; 연속 극한 $(1/4\pi)\int\mathbf n\cdot(\partial_x\mathbf n\times\partial_y\mathbf n)$ \\
    정의 안 됨 & 삼중곱 0, 분모 $\le 0$(공면 120°, 반평행)이면 $Q$를 0이 아닌 ``정의 안 됨''으로 보고 \\
    후보 비교 & 고전 정제, 셀 크기에 맞춘 메시의 LSWT 안정성, 안정할 때만 $E_{\rm cl}+E_{\rm zp}$; 주어진 후보 안의 순위일 뿐 바닥상태 판정이 아님 \\
    검증 & 팔분구 $\pi/2$; 스커미온 $\mp1$(삼각·정사각); 사면체 $|Q|=2$; 삼각 Heisenberg $E=-0.5388J$ \\
    \bottomrule
  \end{tabularx}
  \slidesource{2026-10-01 결정}{contract}
\endgroup


\subsection{결정 기록 (D43): 0.2 공개 API와 옛 API 삭제 일정}\label{decisions-d43-api}
\begingroup% frame-local settings, as in Beamer
  \framesubtitle{2026-10-01 사용자 결정(결정 카드 A: 경고 후 다음 판 삭제).}
  \small\renewcommand{\arraystretch}{1.12}
  \begin{tabularx}{\textwidth}{@{}P{3.0cm}Y@{}}
    \toprule
    항목 & 결정 \\
    \midrule
    공개 API & \texttt{SpinModel}, \texttt{SpinState}, \texttt{solve\_lswt}, \texttt{LSWTSettings}, \texttt{LSWTResult}; README·docstring·예제는 이것만 쓴다 \\
    옛 API & \texttt{SpinSystem}, \texttt{LSWTSolver}, \texttt{SpinOptimizer}, \texttt{EnergyFunction}, SI Hall 메서드, \texttt{lswt} 이름: 0.2에서 동작, 경고에 삭제 판 0.3 명시 \\
    추가 (2026-10-21) & 사용자 선택(``목록에 추가''): \texttt{LSWTHamiltonian}, \texttt{AbstractSolver}, 켈빈 \texttt{Thermodynamics}도 0.2에서 경고, 0.3에서 삭제. 새 API는 이것들을 쓰지 않는다(D65). \texttt{BrillouinZone}은 새 상태 선택 메시가 써서 남긴다 \\
    내부 경로 & 새 API가 내부에서 만드는 \texttt{SpinSystem}은 \texttt{internal\_use()}로 경고하지 않음 \\
    검증 & README 빠른 시작을 \texttt{-W error::DeprecationWarning}으로 실행; $E=-0.5388J$, M점 $\omega=1.0J$ \\
    \bottomrule
  \end{tabularx}
  \slidesource{2026-10-01 결정}{contract}
\endgroup


\subsection{결정 기록 (D44): 공개 배포 전 시각화 묶음}\label{decisions-d44-figures}
\begingroup% frame-local settings, as in Beamer
  \framesubtitle{2026-10-01 Claude 결정(사용자: 공개 배포는 시각화가 갖춰진 뒤), 2026-10-15 사용자 승인(일괄 검토 9번).}
  \small\renewcommand{\arraystretch}{1.12}
  \begin{tabularx}{\textwidth}{@{}P{3.0cm}Y@{}}
    \toprule
    항목 & 결정 \\
    \midrule
    그림 & 중성자 경로·일정 에너지·분말 지도, Berry 곡률과 Chern 수, thermal Hall, 스핀 텍스처와 입체각 밀도, 상도표, 열역학 곡선 \\
    정의 안 됨 & NaN(방향에 무관한 극한이 없는 영모드 가중치 — 자기 Bragg 벡터의 Goldstone 등, D52; 분리 안 된 밴드, 안정 후보 없음, 공면 120°의 $\Omega$)은 회색·빗금·``undefined''; 0이나 이웃 값으로 채우지 않음 \\
    표시 & 저장된 수치만 그림; 단위 변환은 라벨과 함께만; $\langle n\rangle>S$ 구간 음영 \\
    검증 & 120° 상태 M점 $\omega=1.0J$, $Q_x=\pi$ 직선 $\gamma=-1/3$; 벌집 DM Chern $\pm1$; \texttt{code-space/examples/gallery.py} \\
    \bottomrule
  \end{tabularx}
  \slidesource{2026-10-01 결정}{contract}
\endgroup


\subsection{결정 기록 (D45): 표준 LSWT 관측량 그림 보강}\label{decisions-d45-standard}
\begingroup% frame-local settings, as in Beamer
  \framesubtitle{2026-10-01 Claude 결정(사용자 질문에 대한 점검), 2026-10-15 사용자 승인(일괄 검토 10번).}
  \small\renewcommand{\arraystretch}{1.12}
  \begin{tabularx}{\textwidth}{@{}P{3.0cm}Y@{}}
    \toprule
    항목 & 결정 \\
    \midrule
    추가 & 마그논 상태밀도 $g(\omega)$, 성분별 $S^{ab}(\mathbf q,\omega)$, 에너지 적분 세기(확산 + Bragg), $M(h)$와 모멘트 감소 \\
    $M(h)$ & $M=-dE/dh$; 고전 $E_{\rm cl}$과 조화 $E_{\rm cl}+E_{\rm zp}$의 중심 차분; $g(S-\langle n\rangle)$ 근사 쓰지 않음; 불안정 분기 NaN \\
    정의 안 됨 & $\mathbf Q=0$ Bragg 세기(편극인자), 불안정 상태의 $M$과 $\langle n\rangle$은 NaN \\
    검증 & DOS 적분 1; 강자성체 $S^{xx}=S/2$, $S^{zz}=0$; Néel Bragg $m^2/2$; 사각 반강자성체 $M=h/8J$, 포화 위 $M=S$ \\
    \bottomrule
  \end{tabularx}
  \slidesource{2026-10-01 결정}{contract}
\endgroup


\subsection{결정 기록 (D46): 비선형 스핀파 1/S 전개}\label{decisions-d46-nlswt}
\begingroup% frame-local settings, as in Beamer
  \framesubtitle{2026-10-03 Claude 결정(사용자: 비선형 스핀파 solver 선택), 2026-10-15 사용자 잠정 승인(일괄 검토 11번; 엄밀한 검토는 관련 계산을 논문과 비교할 때).}
  \small\renewcommand{\arraystretch}{1.12}
  \begin{tabularx}{\textwidth}{@{}P{3.0cm}Y@{}}
    \toprule
    항목 & 결정 \\
    \midrule
    전개 & LSWT 국소 좌표에서 $s^z=S-n$, $s^+=\sqrt{2S}\,[a-a^\dagger aa/4S]$; 보존 4개까지 자동 생성, $H_n\sim S^{2-n/2}$ \\
    바닥 에너지 & $E_{\rm cl}+E_{\rm zp}+\langle H_4\rangle-6\sum|S^s|^2/(\omega_1+\omega_2+\omega_3)-\tfrac12\zeta^\dagger H(0)^+\zeta$ \\
    마그논 & HF와 tadpole의 정적 보정 + on-shell $\Sigma_3(k,\omega_n)$(붕괴·원천), 확산 $\eta$로 붕괴율 \\
    운동량 & 1/3 칸 이동 메시(합에 닫힘, 영점 없음) 또는 유한 토러스(정확한 RS 섭동) \\
    정의 안 됨 & 비정지 상태, 메시 위 영 모드, $H(0)$ 영 모드 방향 tadpole(질서-무질서), onsite 항 \\
    gap & $\Delta^2=C_\phi/\chi$는 선도 차수 관계; 다음 차수는 무른 모드 2-loop 동역학이 필요 \\
    검증 & 큰 $S$ 정확 대각화(에너지 $10^{-8}$, 마그논 $10^{-5}$); 정사각 $-0.012474J$; 삼각 $-0.0080126J$(문헌 $-0.008029J$와 0.2\%: 문헌 식을 독립 계산하면 $-0.0080126J$로 패키지와 $10^{-8}$ 일치, 차이는 문헌의 $I_3$ 수치 값, 2026-10-13) \\
    \bottomrule
  \end{tabularx}
  \slidesource{2026-10-03 결정}{contract}
\endgroup


\subsection{결정 기록 (D47): 유사 Goldstone gap의 다음 차수}\label{decisions-d47-pseudo-goldstone}
\begingroup% frame-local settings, as in Beamer
  \framesubtitle{2026-10-03 Claude 결정(사용자: ``ㅇ한번 해봐.''), 2026-10-15 사용자 잠정 승인(일괄 검토 12번; 엄밀한 검토는 관련 계산을 논문과 비교할 때).}
  \small\renewcommand{\arraystretch}{1.12}
  \begin{tabularx}{\textwidth}{@{}P{3.0cm}Y@{}}
    \toprule
    항목 & 결정 \\
    \midrule
    극점 & $\det[\omega\sigma_3-H(0)-\Sigma(\omega)]=0$; $H(0)\sim S$, $\Sigma_1\sim S^0$(Nambu 행렬 전체), $\Sigma_2\sim S^{-1}$ \\
    전개 & $\Delta^2=A+B$, $A=K\Sigma_1^{xx}(0)=C_\phi/\chi$; $B$는 $\Sigma_1$의 xx 외 성분과 $\Sigma_2^{xx}(0)$만 필요 \\
    2-loop 정적 & $\Sigma_2^{xx}(0)=U_2-\mathrm{schur}_1$: $\langle x\rangle$ 고정 O($S^0$) 에너지 곡률에서 $\Sigma_1$ 정적 혼합(이중 계산) 제거 \\
    차수 분리 & $H(0)/t$, $\Sigma_1(t\omega)$, $t\Sigma_2$ 축척 후 $t\omega^2=A+Bt$ 적합 \\
    메시 & 1/3 이동 메시를 격자 점군 상으로 평균(C$_3$ 복원) \\
    검증 & 대각 $\Sigma$ $10^{-12}$, Ward $\sim3\times10^{-6}$, Goldstone 극한 $A=B=0$, $A$ vs $C_\phi/\chi$ 0.04\%, $J_{PD}\to0$에서 $B\propto J_{PD}^3$ \\
    NBCP & $S=1/2$: $B/A\approx-27$(Y), $-20$(V) $\Rightarrow$ 1/S 급수가 gap을 정하지 못함 \\
    \bottomrule
  \end{tabularx}
  \slidesource{2026-10-03 결정}{contract}
\endgroup


\subsection{결정 기록 (D48): 패키지 범위는 반고전 툴킷}\label{decisions-d48-scope}
\begingroup% frame-local settings, as in Beamer
  \framesubtitle{2026-10-10 사용자 결정(``추천대로 하자.'').}
  \small\renewcommand{\arraystretch}{1.12}
  \begin{tabularx}{\textwidth}{@{}P{3.0cm}Y@{}}
    \toprule
    항목 & 결정 \\
    \midrule
    범위 & 2D 스핀 모형의 반고전 툴킷: 고전 질서와 그 위의 요동, 중심은 LSWT \\
    안 & 고전 탐색·LT·고전 MC·LL/Langevin 동역학, \texttt{solve\_lswt}, \texttt{solve\_nlswt}(1/S), 관측량·그림 \\
    밖 & 양자 다체 계산(DMRG, QMC, 텐서망); 3D 격자 \\
    ED & 작은 클러스터, 스핀파 결과 검증용으로만 \\
    근거 & LSWT는 고전 바닥 상태가 필요해 한 흐름을 이룸; 양자 다체 도구는 이미 성숙(TeNPy, ITensor, NetKet) \\
    \bottomrule
  \end{tabularx}
  \slidesource{2026-10-10 결정}{contract}
\endgroup


\subsection{결정 기록 (D49): 유한 온도 자화는 $-\partial F/\partial h$}\label{decisions-d49-thermal-magnetization}
\begingroup% frame-local settings, as in Beamer
  \framesubtitle{2026-10-10 Claude 제안(사용자: ``추천대로 진행해''), 2026-10-15 사용자 물리 검토 승인(``승인해'').}
  \small\renewcommand{\arraystretch}{1.12}
  \begin{tabularx}{\textwidth}{@{}P{3.0cm}Y@{}}
    \toprule
    항목 & 결정 \\
    \midrule
    정의 & $M(h,t)=-\partial F/\partial h$, $F=E_{\rm cl}+E_{\rm zp}+t\langle\sum_n\ln(1-e^{-\omega_n/t})\rangle/N_s$; D45의 $M(h)$를 $t>0$으로 잇는다 \\
    API & \texttt{magnetization\_curve(..., temperatures=None)}; 결과에 \texttt{temperatures}, \texttt{thermal} $(n_h, n_t)$를 기본값과 함께 더함(기존 호출 그대로) \\
    계산 & \texttt{harmonic}과 같은 상태 분기에서 $h\pm\delta$ 중앙 차분; 고전 상태는 각 장의 $t=0$ 최소(열적 자기일관성 없음, 저온 근사) \\
    기존 값 & \texttt{ThermalResult.magnetization}$=\sum_i(S_i-\langle n_i\rangle)\hat n_i/N_s$는 그대로 두고, 기울기 각 보정이 없는 모멘트 합이라고 문서에 밝힘 \\
    Goldstone & 반 칸 이동 메시라 영모드가 메시에 없고, 장이 영모드를 0으로 남기면(장 축 둘레 회전, 기울어진 반강자성) $\partial\omega/\partial h\propto k$라 $n_B\,\partial\omega/\partial h$가 유한해 2D에서도 $M$이 유한($\langle n\rangle$은 발산). 2026-10-15 정정: 장이 영모드에 간격을 1차로 열면(강자성 $h=0$, $\partial\omega/\partial h=1$) $n_B\,\partial\omega/\partial h\sim t/\omega$로 2D에서 로그 발산하므로 $t>0$ 값은 NaN과 경고. 판정은 $H(k_0)$의 연한 고윳값 합의 $h$ 기울기 $>10^{-2}$(\texttt{THERMAL\_ZERO\_MODE\_SLOPE}). 그 전에는 격자에 따라 줄어드는 값을 냈다(정사각 강자성 $S=1/2$, $t=0.3$: 격자 12--96에서 0.350$\to$0.151) \\
    불안정 & LSWT가 불안정한 장에서는 NaN \\
    \midrule
    유효 범위 & \texttt{beyond\_lswt} $(n_h,n_t)$와 \texttt{gapless} $(n_h)$: gap이 있으면 열적 $\langle n_i\rangle>S_i$로 표시(\texttt{ThermalResult}와 같은 기준); 영모드·후보가 있으면 $t=0$만 확인하고 $t\ll JS$에서 쓴다고 문서에 밝힘 \\
    기준 상태 & 사용자 결정(2026-10-10): 고전 상태는 $E_{\rm cl}$를 목적함수로 최소화한다($F$가 아님; 차이는 $O(S^0)$로 생략한 4차 항과 같은 차수) \\
    \textbf{열린 항목} & 2D Goldstone에서 유한한 $M$: 사용자 질문(``변경된 정의 때문인가'')에 대한 설명 — 정의 변경의 산물이 아니다. 정확한 이론에서 $-\partial F/\partial h=\langle S^z_{\rm tot}\rangle/N$이고, 발산하는 요동은 장 축 둘레 회전이라 $S^z$를 바꾸지 않는다. 모멘트 합은 그 회전의 종방향 감소만, 기울기 각 항은 기울기 평면 안의 2차 성분을 담아 둘이 상쇄된다. 2026-10-15 검토: $M_{(나)}=M_{(가)}-\sum_iS_i\sin\theta_i\,\delta\theta_i/N$ 유도와 함께 승인. 모멘트 합은 정렬 모멘트 길이 $S_i-\langle n_i\rangle$를 고전 방향에 둔 값으로 남기고 이름은 그대로 둔다 \\
    \bottomrule
  \end{tabularx}
  \slidesource{2026-10-10 결정}{contract}
\endgroup


\subsection{결정 기록 (D50): 약하게 고정된 궤도를 따라 고전 상태를 정제한다}\label{decisions-d50-orbit-refinement}
\begingroup% frame-local settings, as in Beamer
  \framesubtitle{2026-10-10 Claude 제안(사용자: ``추천대로 진행해''), 2026-10-15 사용자 잠정 승인(물리 검토 원칙에 따라 추천안 채택; 엄밀한 검토는 관련 계산에서 논문과 비교할 때; 물리 검토 원칙 절).}
  \small\renewcommand{\arraystretch}{1.12}
  \begin{tabularx}{\textwidth}{@{}P{3.0cm}Y@{}}
    \toprule
    항목 & 결정 \\
    \midrule
    증상 & 로컬(Accelerate BLAS)에서만 \texttt{test\_effective\_potential\_is\_continuous\_with\_d17} 실패; 입력 잡음 seed만 바꿔도 통과·실패가 갈림 \\
    원인 & NBCP Y, $J_{PD}=0.01$, 평면 장 $10^{-3}$: 궤도 곡률 $\sim10^{-9}$ 대 $6.6\times10^{-2}$. L-BFGS-B는 에너지 변화가 반올림 아래라 멈추고, 접선 좌표의 직선 Newton 단계는 궤도를 2차로 벗어나 거부되며, 곡률이 $\sqrt{\epsilon}\max|w|$ 아래면 평탄으로 취급됨. 최소에서 최대 0.3--0.7 rad 떨어진 채 끝남 \\
    결과 & 상태 정확도 추정이 0.2 rad가 되어 \texttt{unsupported\_manifold}; 수렴 못 한 다른 경우는 \texttt{selected}지만 틀린 각($\phi=1.095$, 정답 $\pi/3$) \\
    수정 & \texttt{refine\_classical}의 Newton 뒤 궤도 탐색: 강체 회전에 대한 에너지 미분 $G^{\mathsf T}g$가 $10^3\epsilon\max|w|$보다 크면 그 축을 따라 내리막으로 강체 회전, 각 각도에서 그 축만 빼고 Newton 이완, 첫 부호 변화를 배가 단계로 감싸고 Illinois법으로 근 → 극소. torque와 에너지가 함께 줄 때만 채택 \\
    평탄 궤도 & 정확한 대칭이면 $G^{\mathsf T}g$가 항등적으로 0, 수렴한 상태의 우연 축퇴에서도 반올림 크기라 그대로 둠 \\
    버린 안 & 선택 단계의 정확도 추정에서 궤도 성분을 빼는 안: 판정은 \texttt{selected}가 되지만 수렴 못 한 상태에서 선택 각이 시작 위치에 끌려감(0.747, 1.095) \\
    공유 & 회전 생성자 \texttt{rotation\_generators}를 \texttt{methods.classical}로 옮기고 \texttt{state\_selection}이 씀 \\
    \bottomrule
  \end{tabularx}
  \slidesource{2026-10-10 결정}{contract}
\endgroup


\subsection{결정 기록 (D51): 국소 안정성은 연속 $k$와 영모드 강성으로 판정한다}\label{decisions-d51-local-stability}
\begingroup% frame-local settings, as in Beamer
  \framesubtitle{2026-10-13 Claude 결정(2026-10-01 위임: 추천안대로 진행), 2026-10-15 사용자 잠정 승인(물리 검토 원칙에 따라 추천안 채택; 엄밀한 검토는 관련 계산에서 논문과 비교할 때; 물리 검토 원칙 절).}
  \small\renewcommand{\arraystretch}{1.12}
  \begin{tabularx}{\textwidth}{@{}P{3.0cm}Y@{}}
    \toprule
    항목 & 결정 \\
    \midrule
    문제 & \texttt{solve\_lswt}는 메시 점에서만 $H(k)$를 본다. 메시 사이의 불안정, 특히 Goldstone 모드 옆의 음의 강성($\lambda_{\min}\approx\rho k^2$, $\rho<0$)은 오류 없이 지나간다(NBCP 상 경계 계산 2026-10-13에서 V·Y 상의 장파장 불안정으로 확인) \\
    기준 & 정지 상태의 고전 2차 변분이 모든 $k$에서 $H(k)\succeq0$이면 국소 최소. $H(k)$를 고전 진폭의 2차 형식으로 읽는다(onsite는 D37 coherent-state 값, \texttt{classical\_energy}와 같음) \\
    영모드 밖 & $\lambda_{\min}/\text{scale}$를 이동·비이동 메시($24^2$)의 낮은 점 6개에서 Nelder--Mead로 연속 $k$ 최소화; 영모드 둘레 반지름 0.02(역격자 단위) 공은 뺀다. 0 근처 값은 새 영모드로 등록하고($\pm k_0$) 다시 찾는다 \\
    영모드 안 & $H(k_0)$ 영공간 $V_0$에서 축퇴 섭동 2차까지 해석적으로: 선형 항 $L(u)=V_0^\dagger(u\cdot\partial H)V_0$, 강성 $M(u)=\tfrac12V_0^\dagger(uu:\partial^2H)V_0-V_0^\dagger(u\cdot\partial H)WD^{-1}W^\dagger(u\cdot\partial H)V_0$, $\rho_{\min}=\min_u\lambda_{\min}M(u)$. $H$의 1·2차 도함수가 해석적이라 차분 오차가 없다. 공 안을 반지름 4개×방향 24개로 표본 검사 \\
    판정 & \texttt{unstable}: 음의 고유값, $L\ne0$, $\rho_{\min}<0$ 중 하나. \texttt{marginal}: 음은 없지만 2차가 정하지 못함 — 평탄 강성 방향, $k_0\ne0$ 영모드, $k=0$ 영공간이 깨진 연속 대칭으로 설명되는 수보다 큼(우연 축퇴; NBCP $J_{PD}$·$J_\Gamma$의 3부분격자 궤도, 삼각 하이젠베르크의 장 속 상태), 영모드가 고립되지 않음. \texttt{stable}: 그 밖. \texttt{not stationary}: torque가 허용치보다 크면 판정하지 않음 \\
    API & 새 함수 \texttt{local\_stability(model, state, conditions=None, *, mesh, seeds, ball\_radius, zero\_tolerance, stiffness\_tolerance, stationarity\_tolerance)} $\to$ \texttt{StabilityReport}(최상위 \texttt{stk.local\_stability}); \texttt{QuadraticBoseHamiltonian.second\_derivatives\_at} 추가. 기존 호출은 그대로 \\
    근거 & ``정의되지 않으면 정의되지 않음으로'' 원칙: 우연 축퇴 궤도의 안정성은 고차가 정하므로 \texttt{stable}로 부르지 않고 \texttt{marginal}과 이유를 낸다 \\
    \textbf{열린 항목} & 고차 판정은 D54에서 더함. 자동 호출은 D53에서 정함(\texttt{compare\_states}만) \\
    \bottomrule
  \end{tabularx}
  \slidesource{2026-10-13 결정}{contract}
\endgroup


\subsection{결정 기록 (D52): 영모드 운동량의 한 마그논 가중치는 방향에 무관한 극한이 있을 때만 값으로 낸다}\label{decisions-d52-zero-mode-limit}
\begingroup% frame-local settings, as in Beamer
  \framesubtitle{2026-10-14 사용자 결정(``추천대로 하자''), 2026-10-15 사용자 잠정 승인(물리 검토 원칙에 따라 추천안 채택; 엄밀한 검토는 관련 계산에서 논문과 비교할 때; 물리 검토 원칙 절).}
  \small\renewcommand{\arraystretch}{1.12}
  \begin{tabularx}{\textwidth}{@{}P{3.0cm}Y@{}}
    \toprule
    항목 & 결정 \\
    \midrule
    문제 & $H(q)$에 영모드가 있는 $q$에서 $\mathsf T_q$가 없어 가중치를 모두 NaN으로 냈다. 120° 상태의 핵 Bragg 벡터처럼 가중치가 $q\to G$에서 0으로 가는 곳(보존되는 $S^z$ 총합; A4 닫힌 식)도 NaN이 되어, 정의된 값을 ``정의되지 않음''으로 보고했다 \\
    규칙 & 영모드 운동량 $q_0$에서 $q_0+\delta u$의 가중치를 방향 8개, 반지름 $\delta,\delta/2,\delta/4$($\delta=10^{-4}$×가장 짧은 자기 역격자 벡터)에서 계산. 극한 에너지가 같은 모드는 묶어 합만 쓴다(그 점에서는 합만 정의됨). 모든 방향에서 묶음 가중치가 수렴하고(반지름을 반으로 줄일 때 변화가 줄어듦; $1/\delta$ 발산은 변화가 두 배) 선형 외삽 값이 방향끼리 $10^{-6}$ 안에서 같으면 그 값을 낸다. 아니면 NaN \\
    결과 & 120° 핵 Bragg 벡터·Néel $q=0$: 가중치 0($10^{-13}$). 자기 Bragg 벡터(120° $K$, Néel $(\pi,\pi)$)의 Goldstone 가중치는 발산하므로 NaN 유지. $t>0$에서는 가중치에 Bose 인자가 곱해진다: 가중치가 $|q-q_0|$로 0이 되는 곳(보존 생성원)은 $n_B\sim t/\omega$와 곱해 유한한 극한(고전 $t\chi$; 정사각 Néel $S=1/2$, $t=0.3$: $t/8J=0.0375$와 일치)을 내고, 발산하는 곳은 더 빨리 발산해 NaN. 2026-10-15 정정: 이전 기록의 ``$t>0$에서는 NaN''은 코드 동작과 달랐고 코드가 맞다 \\
    API & \texttt{StructureFactor.limit}(bool 배열) 추가; \texttt{zero\_mode}는 그대로 ``$H(q)$에 영모드가 있음''. 극한인 점의 묶음 가중치는 첫 모드에, 나머지는 0. \texttt{correlation\_path}는 극한이 없는 점만 NaN. 상수 \texttt{ZERO\_MODE\_LIMIT\_*}. 나선 회전틀 \texttt{spiral\_structure\_factor}는 가지($m=0,\pm1$)마다 같은 극한을 쓴다(2026-10-14 추가): 영모드가 있는 가지가 모두 극한을 가지면 \texttt{limit}; 발산하는 가지만 NaN, 다른 가지는 그대로 \\
    근거 & 0을 하드코딩하지 않는다. 극한이 존재하고 방향에 무관할 때만 채우므로 ``정의되지 않으면 정의되지 않음으로'' 원칙과 맞다. 닫힌 식은 테스트에만 쓴다 \\
    \bottomrule
  \end{tabularx}
  \slidesource{2026-10-14 결정}{contract}
\endgroup


\subsection{결정 기록 (D53): \texttt{compare\_states}는 국소 안정성 검사를 기본으로 부르고, \texttt{solve\_lswt}는 부르지 않는다}\label{decisions-d53-compare-stability}
\begingroup% frame-local settings, as in Beamer
  \framesubtitle{2026-10-14 사용자 결정(``추천대로 하자'').}
  \small\renewcommand{\arraystretch}{1.12}
  \begin{tabularx}{\textwidth}{@{}P{3.0cm}Y@{}}
    \toprule
    항목 & 결정 \\
    \midrule
    \texttt{solve\_lswt} & 부르지 않는다. 불안정한 상태의 스펙트럼을 일부러 볼 때가 있고, 연속 $k$ 최소화는 비용이 든다. 메시 통과가 안정성 검사가 아니라는 점은 docstring이 \texttt{local\_stability}를 안내한다 \\
    \texttt{compare\_states} & 새 인자 \texttt{local\_stability=True}(기본). 후보마다 D51 판정을 \texttt{CandidateReport.stability}와 \texttt{to\_dict()["local\_stability"]}에 붙인다. 메시는 통과했지만 판정이 \texttt{unstable}이면(메시 사이·Goldstone 옆 음의 강성) \texttt{status="unstable"}, 영점 에너지 NaN: 불안정한 2차 배경에는 조화 진공이 없다(D42와 같은 이유). \texttt{marginal}·\texttt{not stationary}는 붙이기만 하고 상태를 바꾸지 않는다 \\
    근거 & \texttt{compare\_states}는 바닥 상태 후보를 고르는 함수라, 메시가 놓친 불안정한 후보를 안정으로 순위에 올리면 물리적으로 틀린 답이 된다. 끄면(\texttt{False}) 예전 메시 판정 \\
    시험 & 정사각 강자성 $J_3=0.2505$: 메시만으로는 \texttt{stable}, 기본값으로는 \texttt{unstable}; 삼각 120°는 \texttt{stable} \\
    \bottomrule
  \end{tabularx}
  \slidesource{2026-10-14 결정}{contract}
\endgroup


\subsection{결정 기록 (D54): 2차가 정하지 못한 영모드는 고차 항으로 판정한다}\label{decisions-d54-higher-order-stability}
\begingroup% frame-local settings, as in Beamer
  \framesubtitle{2026-10-14 Claude 결정(사용자 ``1,2번은 진행하자''), 2026-10-15 사용자 잠정 승인(물리 검토 원칙에 따라 추천안 채택; 엄밀한 검토는 관련 계산에서 논문과 비교할 때; 물리 검토 원칙 절).}
  \small\renewcommand{\arraystretch}{1.12}
  \begin{tabularx}{\textwidth}{@{}P{3.0cm}Y@{}}
    \toprule
    항목 & 결정 \\
    \midrule
    문제 & D51은 평탄 강성 방향, $k_0\ne0$ 영모드, 우연 축퇴를 모두 \texttt{marginal}로 둔다. 고차에서 올라가는 상태(국소 최소), 내려가는 상태(안장점), 정확히 평평한 고전 축퇴 다양체가 구분되지 않는다 \\
    평탄 강성 & $H(k)$는 $k$에 대해 정확하므로 고차 항도 조화 진폭의 문제다. $\lambda_{\min}(H(k_0+q))$를 $q$ 방향 360개(가장 낮은 곳 근처 다듬기)에서 최소로 하여 $|q|=r,r/2,\dots,r/16$($r$=D51 공 반지름)에서 잰다 \\
    우연 영모드 & 진폭의 문제. $k_0$를 $k=0$으로 접는 초격자(분모 6 이하, 스핀 144개 이하)에서 고전 헤시안 영공간에서 에너지를 바꾸지 않는 전역 회전의 생성자를 뺀 부분공간을 잡는다. 그 방향(1차원: 1개, 2차원: 8개, 6차원까지: 기저와 쌍)과 반대 방향으로 스핀당 rms 각 $0.2,0.1,\dots,0.0125$ rad만큼 움직이고, 나머지 단단한 방향은 L-BFGS로 이완한 스핀당 에너지 변화를 잰다. 이완 덕에 4차 계수는 단단한 모드와의 3차 결합을 포함한 유효값이다 \\
    읽는 법 & 바닥값(scale의 $10^{-11}$, 운동량은 \texttt{zero\_tolerance})을 넘는 가장 작은 진폭에서 부호, 그 위 진폭과의 비로 차수 $p=\log_2[f(2a)/f(a)]$. 어느 방향에서든 내려가면(홀수 3차 포함) \texttt{unstable}; 모든 방향에서 올라가면 그 영모드는 해가 없음(다른 문제가 없으면 \texttt{stable}); 어떤 방향에서 가장 큰 진폭까지 바닥값 아래면 고전 축퇴 다양체로 \texttt{marginal}(요동이 정함: order by disorder); 두 연속 진폭에서 같은 부호로 $p\ge2.5$가 안 나오면 미해결로 \texttt{marginal} \\
    검산 & 쉬운 평면 강자성 임계장 $h=4S\delta$: $S^2\delta\theta^4/2$와 $10^{-3}$. 정사각 반강자성 포화장: $(\pi,\pi)$ 영모드를 4스핀 셀로 접어 $JS^2\theta^4$. $S=1$ 쉬운 축 단일 이온 Néel의 연화장 $h=0.9$: 4차로 내려감(1차 spin flop). XXZ 반강자성 $h=4S\sqrt{J_z^2-J_\perp^2}$, 장 속 삼각 하이젠베르크, NBCP Y($J_{PD}=0.01$): 평평(고전 축퇴). 정사각 강자성 $J_3=1/4$: $q^4$로 올라감 \\
    API & \texttt{local\_stability(..., higher\_order=True)}(끄면 D51 판정); \texttt{StabilityReport.higher\_order}(\texttt{HigherOrderCheck}: kind, fractional, spins, directions, outcome, order, lowest, amplitude). 판정 이름은 그대로. 상수 \texttt{STABILITY\_AMPLITUDE}·\texttt{\_FLOOR}·\texttt{\_FOLD\_DENOMINATOR}·\texttt{\_FOLD\_SPINS}. \texttt{compare\_states}는 D53대로 \texttt{unstable}만 상태를 바꾸므로 고차에서 내려가는 후보도 불안정으로 뺀다 \\
    근거 & ``정의되지 않으면 정의되지 않음으로'': 고차가 정하는 것은 정하고, 고전적으로 정해지지 않는 것(축퇴 다양체)은 \texttt{marginal}로 남기되 이유를 ``요동이 정함''으로 구체화한다. ``평평''은 바닥값 안에서 평평하다는 뜻이지 정확함을 증명한 것이 아니다 \\
    한계 & 부분공간 방향을 표본으로만 보므로 표본 사이에만 있는 내리막은 놓칠 수 있다. 영모드 선(고립되지 않은 영모드)과 비정합 $k_0$는 고차 판정하지 않는다 \\
    \bottomrule
  \end{tabularx}
  \slidesource{2026-10-14 결정}{contract}
\endgroup


\subsection{결정 기록 (D55): 단일 $Q$ 나선의 국소 안정성은 회전틀 $H(k)$로 판정한다}\label{decisions-d55-spiral-stability}
\begingroup% frame-local settings, as in Beamer
  \framesubtitle{2026-10-14 Claude 결정(사용자 ``다음 개발항목을 정하자'', 추천안), 2026-10-15 사용자 잠정 승인(물리 검토 원칙에 따라 추천안 채택; 엄밀한 검토는 관련 계산에서 논문과 비교할 때; 물리 검토 원칙 절).}
  \small\renewcommand{\arraystretch}{1.12}
  \begin{tabularx}{\textwidth}{@{}P{3.0cm}Y@{}}
    \toprule
    항목 & 결정 \\
    \midrule
    문제 & D51은 정합 \texttt{SpinState}만 받았다. 비정합·긴 주기 나선은 큰 초격자를 만들어야 했다(NBCP 상 경계: 36–435 스핀) \\
    방법 & \texttt{local\_stability}가 \texttt{IncommensurateStructure}를 받으면 D34의 회전틀 모형(\texttt{rotating\_frame\_model})에 D51·D54를 그대로 적용. 변수 변환 $S=R_n(\phi)S'$은 모든 배치에서 정확하므로 회전틀 $H(k)$가 실험실 2차 변분이고, 어떤 $Q$든 원시 셀 하나로 연속 $k$를 본다 \\
    Goldstone & $n$ 둘레 회전: 회전틀 $k=0$. $n$에 수직인 축의 회전도 대칭이면(등방 교환, 장 없음) 칸 $R$에서 축 $R_n(-Q\cdot R)m$ 둘레 회전이라 회전틀 $k=\pm Q$에 영모드 하나씩: nullity 1이면 \texttt{goldstone}. D54 진폭 검사에서도 이 생성자를 뺀다 \\
    피치 & 브라베 격자의 평면 나선은 어떤 피치에서도 torque가 0. 수직 회전이 대칭이면 $\partial E/\partial Q\ne0$이 $\pm Q$의 선형 항으로 나타나 \texttt{unstable}. 수직 회전이 대칭이 아니면 틀린 피치는 꼬인 상태로 국소 안정일 수 있다(XY의 균일한 꼬임처럼). 그래서 \texttt{settings["pitch\_gradient"]}에 $\partial E/\partial Q$를 두고 0이 아니면 요약에 ``twisted spiral''로 적는다. 판정은 국소 안정성만 말한다 \\
    검산 & 삼각 120° 나선: \texttt{stable}, 영모드 $0,\pm K$ 모두 Goldstone, 초격자 판정과 같음. $J_1$–$J_2$ 사슬 비정합 $Q$: \texttt{stable}(영모드 3개); 틀린 피치: 선형 항으로 \texttt{unstable}. 쉬운 평면 $J_z/J=0.9$: Goldstone 하나로 \texttt{stable}, 피치 $+0.02$는 꼬인 상태로 \texttt{stable}. 쉬운 축 1.1: \texttt{unstable}. 축 방향 장 속 원뿔: \texttt{stable} \\
    API & \texttt{local\_stability(model, state, ...)}의 \texttt{state}가 \texttt{SpinState} 또는 \texttt{IncommensurateStructure}. 보고서의 운동량은 회전틀 원시 역격자 좌표. $n$ 둘레 U(1) 대칭이 없는 모형은 D34대로 \texttt{SpiralSymmetryError} \\
    한계 & 비정합 $k_0$의 우연 영모드는 셀로 접을 수 없어 고차 판정하지 않는다(\texttt{marginal}) \\
    \bottomrule
  \end{tabularx}
  \slidesource{2026-10-14 결정}{contract}
\endgroup


\subsection{결정 기록 (D56): 종방향 두 마그논 연속체는 상호작용 없는 $O(S^0)$ 항으로 따로 낸다}\label{decisions-d56-two-magnon}
\begingroup% frame-local settings, as in Beamer
  \framesubtitle{2026-10-14 사용자 선택(``두 마그논 연속체''), 구현은 Claude 결정, 2026-10-15 사용자 잠정 승인(물리 검토 원칙에 따라 추천안 채택; 엄밀한 검토는 관련 계산에서 논문과 비교할 때; 물리 검토 원칙 절).}
  \small\renewcommand{\arraystretch}{1.12}
  \begin{tabularx}{\textwidth}{@{}P{3.0cm}Y@{}}
    \toprule
    항목 & 결정 \\
    \midrule
    문제 & 구조인자는 한 마그논($O(S)$)과 탄성만 낸다. 중성자 실험의 연속체 세기 중 종방향 두 마그논($O(S^0)$)이 빠져 있다 \\
    식 & $\delta S_i^a\supset-n_i^a a_i^\dagger a_i$. $L^b(q)^\dagger$의 두 생성 부분 $C^b_{nn'}(k)=-N^{-1/2}\sum_i n_i^b T_{N_s+i,N_s+n}(-k-q)T_{i,N_s+n'}(k)$, 대칭화 $A=C_{nn'}(k)+C_{n'n}(-k-q)$, $S_2^{ab}(q,\omega)=\tfrac12\sum_{k,nn'}A^{a*}A^b\,\delta(\omega-E_{N_s+n}(-k-q)-E_{N_s+n'}(k))$. 에너지는 $T$의 hole 열 고윳값으로 읽어 띠 이름표 대응을 가정하지 않는다 \\
    범위 & 상호작용 없는 연속체만. 같은 차수의 한 마그논 에너지·가중치 1/S 보정, 다음 HP 항의 가로 가중치 보정, 비공선 상태의 마그논 붕괴는 포함하지 않으며 docstring과 노트에 밝힌다. $t=0$만($t>0$의 열 마그논 산란 $b^\dagger b$는 \texttt{NotImplementedError}). 실험실 틀 결과만 \\
    영모드 & 내부 운동량 메시는 반 칸 이동 메시. $k$ 또는 $-k-q$에 영모드가 있는 점은 빼고 \texttt{excluded}에 센다(측도 0, 가중치는 적분 가능) \\
    검산 & 같은 토러스에서 실공간 Wick 수축 $\langle\delta n\,\delta n\rangle_c=\langle a^\dagger a\rangle\langle aa^\dagger\rangle+\langle a^\dagger a^\dagger\rangle\langle aa\rangle$과 정적 합이 $10^{-16}$에서 일치(정사각 Néel, 삼각 120°). 장 속 강자성: 0. Néel 쌍 에너지 $=\omega(k)+\omega(k+q)$, 가중치는 엇갈림 축 $zz$만. 스펙트럼 적분 = 정적 합 \\
    API & 새 함수 \texttt{observables.two\_magnon.two\_magnon\_continuum(result, q\_points, mesh=(24, 24), temperature=0.0)} $\to$ \texttt{TwoMagnonContinuum}(\texttt{energies}, \texttt{weights}, \texttt{static()}, \texttt{spectrum()}); 한 마그논 \texttt{StructureFactor}와 합치지 않고 따로 둔다(차수와 범위가 다르므로) \\
    \bottomrule
  \end{tabularx}
  \slidesource{2026-10-14 결정}{contract}
\endgroup


\subsection{결정 기록 (D57): 마그논 붕괴율은 넓힘 없이 황금률을 삼각형 적분으로 계산한다}\label{decisions-d57-decay}
\begingroup% frame-local settings, as in Beamer
  \framesubtitle{2026-10-14 사용자 선택(``마그논 붕괴''), 구현은 Claude 결정, 2026-10-15 사용자 잠정 승인(물리 검토 원칙에 따라 추천안 채택; 엄밀한 검토는 관련 계산에서 논문과 비교할 때; 물리 검토 원칙 절).}
  \small\renewcommand{\arraystretch}{1.12}
  \begin{tabularx}{\textwidth}{@{}P{3.0cm}Y@{}}
    \toprule
    항목 & 결정 \\
    \midrule
    문제 & D46의 붕괴율 $-\mathrm{Im}\,\Sigma_3(k,\omega_n(k)+i\eta)$은 메시 위 극점의 합이라 인위적 넓힘 $\eta$와 메시에 함께 의존한다. 물리량인 $\eta\to0$ 극한이 없었다 \\
    식 & $\Gamma_n(k)=\pi\langle\sum_{n_1n_2}2|D_{n_1n_2;n}(q,k-q)|^2\,\delta(\omega_n(k)-\omega_{n_1}(q)-\omega_{n_2}(k-q))\rangle_q$. D46과 같은 대칭화 꼭짓점 $D$(대형 $S$ ED로 검증됨), LSWT 에너지로 계산하는 on-shell Born 근사, 상대 차수 $1/S$ \\
    적분 & 자기 브릴루앙 영역을 $N_1\times N_2$ 격자(1/3칸 이동)의 삼각형으로 나누고, 에너지 차와 $2|D|^2$를 삼각형마다 선형 보간해 델타 함수를 정확히 적분한다(3차원 사면체법의 2차원판). 연속체 아래에서는 정확히 0 \\
    문턱 & \texttt{two\_magnon\_threshold}: 두 마그논 연속체 아래 끝 $\min_{q,n_1,n_2}\omega_{n_1}(q)+\omega_{n_2}(k-q)$를 격자 최소와 Nelder--Mead 정밀화로 구한다. 운동학 조건일 뿐 충분조건이 아니다: 병진$\times$스핀 회전 대칭이 펼친 운동량을 보존하면 열린 통로의 꼭짓점이 0이다(기울어진 정사각 반강자성의 접힌 위 띠) \\
    한계 & 연속체 끝 근처에서는 자기에너지가 특이해져 Born 근사가 깨진다(Zhitomirsky--Chernyshev 2013). 격자 한 칸보다 좁은 붕괴 영역은 0으로 나오므로 격자를 키워야 한다. 자기무모순 극 방정식은 풀지 않는다(더 높은 차수를 섞으므로) \\
    검산 & 삼각형 적분 $=$ 해석적 상태밀도. 기울어진 정사각 반강자성 $h=0.9h_c$에서 $\eta=0.1$, $24\times24$ 메시의 $-\mathrm{Im}\,\Sigma_3$와 10\% 안에서 일치. 같은 상태에서 Goldstone 가지가 대각선 방향으로 볼록해지는 장 $0.755929\,h_c$($2/\sqrt7$과 $10^{-6}$ 안에서 같음, 수치로 찾음)는 Zhitomirsky--Chernyshev 1999의 $h^*\approx0.76h_c$와 맞고, $0.7$, $0.74h_c$에서 $\Gamma=0$, $0.9h_c$에서 $\Gamma>0$. 3차 항이 없는 Néel: 0. 삼각 120°: 연속체 아래 띠는 0, $48\to96$ 격자에서 5\% 안 \\
    API & \texttt{NLSWTResult.decay\_rates(k\_points, grid=(64, 64))} $\to$ \texttt{\{lswt, decay\_rate\}}; \texttt{NonlinearSpinWaves.decay\_rate}; \texttt{observables.two\_magnon.two\_magnon\_threshold(result, k\_points, grid=(48, 48))}. 기존 \texttt{magnon\_energies(..., broadening)}은 그대로 둔다 \\
    \bottomrule
  \end{tabularx}
  \slidesource{2026-10-14 결정}{contract}
\endgroup

\subsection{결정 기록 (D58): 마그논 스펙트럼 함수는 on-shell을 기본으로, Dyson 재합은 진단값과 함께 낸다}\label{decisions-d58-spectral}
\begingroup% frame-local settings, as in Beamer
  \framesubtitle{2026-10-14 사용자 선택(``상호작용 스펙트럼''), 구현은 Claude 결정, 2026-10-15 사용자 잠정 승인(물리 검토 원칙에 따라 추천안 채택; 엄밀한 검토는 관련 계산에서 논문과 비교할 때; 물리 검토 원칙 절).}
  \small\renewcommand{\arraystretch}{1.12}
  \begin{tabularx}{\textwidth}{@{}P{3.0cm}Y@{}}
    \toprule
    항목 & 결정 \\
    \midrule
    문제 & 1/S 에너지 이동(D46), 붕괴율(D57), 두 마그논 연속체(D56)가 따로 있어 하나의 $S(q,\omega)$로 볼 수 없었다 \\
    자기에너지 & $\Pi^{(3)}_n(k,\omega)=\int\rho^{\rm dec}(x)/(\omega-x+i0)\,dx-\int\rho^{\rm src}(x)/(\omega+x+i0)\,dx$. 두 밀도는 D57 삼각형 적분으로 에너지 격자에서 구하고, 힐베르트 변환은 조각별 선형 보간에 대해 정확히 계산한다(자기에너지 안에 넓힘 없음) \\
    기본(on-shell) & 엄밀한 상대 $1/S$ 차수: $\tilde\varepsilon=\varepsilon+\delta\varepsilon+{\rm Re}\,\Pi^{(3)}(\varepsilon)$에 반폭 $\Gamma+\eta$인 로렌츠 함수 \\
    Dyson(선택) & $A=-\frac1\pi{\rm Im}[\omega+i\eta-\varepsilon-\delta\varepsilon-\Pi^{(3)}(\omega)]^{-1}$. 선 모양, 연속체로 옮겨 간 무게, 연속체 끝에 붙는 봉우리를 주지만 $1/S$로 통제되지 않는다. 진단값: $\omega<0$ 무게는 보손에서 $\le0$이어야 하므로 양수면 가짜 극이다. 삼각 120° $S=5$: 무게 합 1(오차 $10^{-2}$), 음주파수 무게 음수. $S=1/2$: 연속체 끝의 접힌 띠에서 정적 이동과 ${\rm Re}\,\Pi^{(3)}$가 LSWT 에너지보다 커 양의 무게를 가진 가짜 극이 생기고, 진단값이 이를 드러낸다 \\
    띠 처리 & 띠 대각만 쓴다(비대각은 더 높은 차수). 정확히 겹친 띠 묶음에서는 대각 성분이 기저에 의존하므로 묶음 평균(대각합)을 쓴다 \\
    $S(q,\omega)$ & \texttt{dressed\_structure\_factor}: $\sum_nW_n(q)A_n(q,\omega)+S_2(q,\omega)$. 완전한 $O(S^0)$가 아니다: 한 마그논 가중치의 $1/S$ 보정과 비공선 상태의 한·두 마그논 간섭 항은 같은 차수지만 넣지 않으며 docstring과 노트에 밝힌다. $t=0$, 실험실 틀, 영모드 운동량은 거부 \\
    검산 & 반원 밀도의 힐베르트 변환 해석해($2\times10^{-3}$), on-shell 붕괴율 = D57(1\%), 3차 항 없는 Néel은 이동된 로렌츠 함수. 연속체 끝의 띠에서 ${\rm Re}\,\Pi^{(3)}$: 삼각형 격자 48/96/160 $\to-1.656/-1.649/-1.653$, 기존 메시 합 12--48 $\to-1.597\dots-1.661$(끝의 특이성 때문에 메시 합이 느리게 다가간다) \\
    API & \texttt{NLSWTResult.spectral\_function(k\_points, omega, eta, mode="on\_shell", grid, levels)}; \texttt{NonlinearSpinWaves.cubic\_spectral\_densities}; \texttt{engine.triangle\_density}, \texttt{engine.hilbert\_transform}; \texttt{observables.interacting\_spectrum.dressed\_structure\_factor(lswt\_result, nlswt\_result, q\_points, omega, eta, component, mode)} \\
    \bottomrule
  \end{tabularx}
  \slidesource{2026-10-14 결정}{contract}
\endgroup

\subsection{결정 기록 (D59): 동적 구조인자를 $O(S^0)$까지 빠짐없이 계산한다}\label{decisions-d59-structure-s0}
\begingroup% frame-local settings, as in Beamer
  \framesubtitle{2026-10-14 사용자 선택(``완전한 S⁰ 스펙트럼''), 구현은 Claude 결정, 2026-10-15 사용자 잠정 승인(물리 검토 원칙에 따라 추천안 채택; 엄밀한 검토는 관련 계산에서 논문과 비교할 때; 물리 검토 원칙 절).}
  \small\renewcommand{\arraystretch}{1.12}
  \begin{tabularx}{\textwidth}{@{}P{3.0cm}Y@{}}
    \toprule
    항목 & 결정 \\
    \midrule
    문제 & D58의 $\sum_nW_nA_n+S_2$에는 같은 차수의 두 항이 빠져 있었다: 한 마그논 가중치의 $1/S$ 보정과, 비공선 상태에서 3차 꼭짓점이 섞는 한·두 마그논 간섭 \\
    스핀 연산자 & 해밀토니안과 같은 HP 절단과 같은 응축(D46)으로 $S^a_q=(\mathsf A+\delta\mathsf A)\psi_q+L^a_q$. $\delta\mathsf A$: $s^+$의 $-a^\dagger aa/4S$의 HF 축약($2\langle a^\dagger a\rangle a+\langle aa\rangle a^\dagger$)과 $-n\,a^\dagger a$ 속 응축 $\alpha$(정렬 모멘트의 $1/S$ 회전) \\
    상관함수 & $\chi=\mathsf ag\mathsf a^\dagger+\delta\mathsf a g\mathsf a^\dagger+{\rm h.c.}+\mathsf agMg\mathsf a^\dagger+\mathsf ag\mathsf T^\dagger\sigma_3K+{\rm h.c.}+\chi_{LL}$. $M$은 HF·tadpole·3차 자기에너지의 남부 전체 행렬, $K=\langle\langle j;L^\dagger\rangle\rangle$은 운동방정식 $\langle\langle\psi;L^\dagger\rangle\rangle=G_0\sigma_3\langle\langle j;L^\dagger\rangle\rangle$에서. $\Pi,K,\tilde K,\chi_{LL}$은 연산자 $(j,L)$의 자유 두 마그논 거품 하나의 블록 \\
    최종 상태 & 한 마그논: 극 $\varepsilon+d_j$($d_j$는 $M_{gg}(\varepsilon)$ 고유값), 잔여 = LSWT 가중치 + 준입자 잔여($M'$) + 연산자 보정 + 띠간·비정상 자기에너지에 의한 고유벡터 회전 + 간섭. 두 마그논: $\tfrac12|c_L+\mathsf ag(E_P)\mathsf T^\dagger\sigma_3X|^2$(종방향 + 가상 마그논을 거친 가로 성분; 교차항이 간섭, 제곱이 연속체 속 꼬리) \\
    열역학 극한 & 거품은 삼각형 적분(D57) 행렬 밀도의 코시 변환. 극 부분을 빼면 유한한 나머지가 남고, 이를 넓힘 0·양의 주파수에서 계산해(두 마그논 상태만 기여, 음주파수 구조 배제) $\eta$ 로렌츠 함수로 다시 변환한다. 극은 $R_j/(z-\varepsilon-d_j)$로 재합: 붕괴하는 마그논은 ${\rm Im}\,d=-\Gamma$(D57)이고 잔여가 복소수(Fano 비대칭). 그 밖에서는 엄밀한 $O(S^0)$ \\
    겹친 띠 & 정확히 겹친 묶음은 $M_{gg}$를 대각화해 투영자 $P_j$로 나눈다. 대칭으로 생긴 겹침에서는 $M'$과 $P_j$가 교환해 정확하다 \\
    범위 & 쌍 에너지는 LSWT 값(두 마그논 상호작용·속박 상태는 더 높은 차수). $t=0$, 실험실 틀. $O(S^0)$ 가중치가 LSWT와 비슷해지면(예: $S=1/2$ 붕괴 문턱 근처) 통제되지 않으며, 음의 잔여나 넓힘 너머의 음의 스펙트럼이 그 신호다. 붕괴하는 마그논의 복소 잔여는 봉우리에서 먼 곳에 천천히 줄어드는 분산형 꼬리를 남긴다(on-shell 인공물) \\
    검산 & 두 스핀 클러스터(3차·tadpole 있음)의 큰 $S$ 정확 대각화($S=10\dots31$): 한 마그논 $O(S^0)$ 가중치 4개가 $3\times10^{-4}$, 두 마그논 상태 6개가 $10^{-4}$ 안에서 일치($q=0$과 역격자 벡터). Néel 합 규칙 $\sum\langle S^aS^a\rangle=S(S+1)$의 $S^0$ 항 $10^{-12}$. 강자성 보정 0. Néel 연속체 = D56(3\%). 붕괴 없는 기울어진 반강자성 $h=0.5h_c$: 열역학 극한과 메시 합이 2\% 안. 삼각 120° $S=5$: 극 위치·폭 = D58 on-shell, 가중치 10--15\% 감소. 붕괴하는 띠($\varepsilon=11.75$, $\Gamma\approx0.078$): $\varepsilon+i0$의 엄밀한 잔여는 격자 $36\dots144$에서 LSWT의 $0.71\dots0.92$로 수렴하지 않고, 선폭 평균 무게는 격자 $48\dots96$에서 $0.886\dots0.908$. 창 $10.5<\omega<13$의 전체 스펙트럼 적분은 3\% 안 \\
    붕괴하는 띠의 무게 & $\mathsf M'(\varepsilon+i0)$은 두 마그논 밀도의 van Hove 구조를 따라 선보다 날카롭다. 보고하는 무게와 한 마그논 선은 $\mathsf M'$, 비대각 $\mathsf M$, $\mathsf K$를 $\varepsilon+i\Gamma$에서 계산(선폭 평균)하고, 위치와 폭은 $\mathsf M_{gg}(\varepsilon+i0)$. 차이는 $\Gamma\,\partial^2\Pi$ 차수로 $1/S$ 정확도 밖. 전체는 엄밀한 식, 연속체 = 전체 $-$ 선이라 붕괴하는 선 근처 몇 $\Gamma$ 안의 선·연속체 나눔은 관례이고, 그 창의 연속체 모양은 격자에 따라 선 무게의 몇 \% 흔들린다 \\
    API & \texttt{observables.interacting\_spectrum.structure\_factor\_s0(nlswt\_result, q\_points, omega, eta, component, grid, levels)}(연속 극한), \texttt{structure\_factor\_states(nlswt\_result, q\_points)}(유한 토러스 최종 상태); \texttt{methods.nlswt.structure}; \texttt{engine.cauchy\_transform}, 여러 함수용 \texttt{triangle\_density}, \texttt{NonlinearSpinWaves.bilinear\_pair\_amplitudes} \\
    \bottomrule
  \end{tabularx}
  \slidesource{2026-10-14 결정}{contract}
\endgroup

\subsection{결정 기록 (D60): 두 마그논 연속체를 유한 온도로 넓힌다}\label{decisions-d60-two-magnon-thermal}
\begingroup% frame-local settings, as in Beamer
  \framesubtitle{2026-10-14 개발 순서에 따른 항목(사용자 "순서대로 알아서 진행"), 구현은 Claude 결정, 2026-10-15 사용자 잠정 승인(물리 검토 원칙에 따라 추천안 채택; 엄밀한 검토는 관련 계산에서 논문과 비교할 때; 물리 검토 원칙 절).}
  \small\renewcommand{\arraystretch}{1.12}
  \begin{tabularx}{\textwidth}{@{}P{3.0cm}Y@{}}
    \toprule
    항목 & 결정 \\
    \midrule
    문제 & D56의 종방향 연속체는 $t=0$만 다뤘고 $t>0$은 거부했다. 유한 온도에서는 열 마그논의 산란($b^\dagger b$)과 쌍 소멸이 더해진다 \\
    식 & $\hat L^\beta(\mathbf q)^\dagger$를 $\mathsf T$의 모든 $2N_s$ 열로 쓰고 자유 보손의 열평형 Wick 정리: $S_2=\tfrac12\sum A^*A\,f_\sigma f_\tau\,\delta(\omega-e_\sigma-e_\tau)$, 생성 연산자 $e=+E$, $f=1+n_B$, 소멸 연산자 $e=-E$, $f=n_B$. 쌍 생성 $(1+n_1)(1+n_2)$, 산란 $\pm(\varepsilon_1-\varepsilon_2)$에서 $(1+n_1)n_2$, 쌍 소멸 $n_1n_2$ \\
    온도 단위 & \texttt{thermal\_quantities}와 같이 $t=k_BT$를 모델 에너지 단위로 \\
    영모드 & $t>0$에서 영모드 근처 산란 무게가 $n_B|A|^2$로 커져 2차원 영역 합이 로그 발산(Mermin--Wagner, 보손 수와 같음). 따라서 영모드 판정은 D25와 같은 방식(\texttt{scan\_zero\_modes}, 후보가 있으면 \texttt{gapless} 지정 요구)이고, 간격 없는 스펙트럼이면 무게를 NaN으로 보고하고 경고한다. 메시 값으로 채우지 않는다 \\
    범위 & 상호작용 없는 연속체. 고전 상태 고정, 1/S 보정 없음. 연결 상관만이라 Bragg 벡터의 열적 모멘트 감소(탄성 부분)는 포함하지 않는다. D59의 상호작용 구조인자는 여전히 $t=0$ \\
    검산 & 쉬운 축 Néel($J_z/J=1.05$, 장 $0.05$, 간격 있음): 등시간 합 = 열 수축의 실공간 Wick 계산 $10^{-13}$($t=0.3,1$). 상세 균형 $\sum W(-\mathbf q)=\sum W(\mathbf q)e^{-e/t}$가 $10^{-10}$. $t\to0$ = D56. 장 속 강자성: 산란만, $\langle n_k(1+n_{k+q})\rangle$와 $10^{-12}$. 간격 없는 Néel은 NaN \\
    API & \texttt{two\_magnon\_continuum(result, q\_points, mesh, temperature, zero\_modes, gapless)}; 결과 \texttt{settings}에 \texttt{gapless}, \texttt{gapless\_decision} \\
    \bottomrule
  \end{tabularx}
  \slidesource{2026-10-14 결정}{contract}
\endgroup

\subsection{결정 기록 (D61): 물리적이지 않거나 정의되지 않는 값은 NaN으로 낸다}\label{decisions-d61-undefined-nan}
\begingroup% frame-local settings, as in Beamer
  \framesubtitle{2026-10-15 사용자 지시(``물리적인 상황이 아닌 경우, 또는 정의할 수 없는 상황에서는 숫자를 주지말고, nan으로 해''), 적용 범위는 Claude 결정, 사용자 확인 대기.}
  \small\renewcommand{\arraystretch}{1.12}
  \begin{tabularx}{\textwidth}{@{}P{3.0cm}Y@{}}
    \toprule
    항목 & 결정 \\
    \midrule
    원칙 & 입력이 물리적이지 않거나(불안정, 비정지, $\langle n_i\rangle>S_i$) 극한이 없으면(발산, 방향·메시 의존) 숫자 대신 NaN과 경고를 낸다. 0은 유도된 물리값일 때만 쓴다(보존 생성자의 가중치, $T=0$의 열 항 등) \\
    불안정 상태 & \texttt{solve\_lswt}: 정규화(\texttt{MAGSWT}, \texttt{k-dependent})는 영모드만 들어 올린다. $H(k)$의 고윳값이 $-\texttt{zero\_mode\_tolerance}\cdot\max|\mathrm{eig}|$ 아래면 정규화를 골라도 \texttt{LSWTError}. 이전에는 MAGSWT 이동량에 상한이 없어 불안정 상태에도 에너지를 냈다 \\
    비정지 상태 & 토크가 남으면 보손 선형항이 남아 조화 차수가 없다. 영점·바닥 에너지와 보손 수를 NaN으로(이전에는 경고만). $H(k)$와 그 고윳값은 이차 부분으로 남긴다 \\
    고정 궤도 영점 에너지 & \texttt{LSWTZeroPointEnergy}(constrained): 투영 뒤 $H(k)$가 양의 준정부호가 아니면 NaN. 이전에는 복소 쌍만 불안정으로 세고 실수 진동수의 음에너지 모드(Landau 불안정)를 양의 진동수로 잘못 셌다 \\
    LSWT 밖 온도 & \texttt{thermal\_quantities}: \texttt{beyond\_lswt}인 온도에서 $F,U,S,C$, 모멘트, 모멘트 합을 NaN으로(보손 수는 원인을 보이려 남긴다). \texttt{ordered\_moments}는 $\langle n_i\rangle>S_i$에서 NaN. \texttt{magnetization\_curve}도 같은 곳에서 NaN \\
    열적 Bragg & 구조인자의 탄성 부분이 $t>0$에서도 $T=0$ 모멘트를 썼다. 열 모멘트 $S_i-\langle n_i\rangle(t)$를 쓰고, 간격 없는 스펙트럼에서는 NaN(2차원에서 질서 없음) \\
    1/S & \texttt{magnon\_energies}: 음수인 1/S 에너지는 NaN. 열역학 극한에서 $\eta=0$의 \texttt{decay\_rate}는 NaN(\texttt{decay\_rates} 사용). Dyson 진단 \texttt{negative\_frequency\_weight}는 음의 진동수 격자가 없으면 NaN(확인 안 함) \\
    기타 & MC 잭나이프 오차: 블록이 둘 미만이거나 측정 수보다 많으면 NaN. 두 마그논 연속체에서 영모드로 뺀 쌍의 에너지 NaN(가중치 0). 옛 \texttt{JH\_method}가 복소 고윳값 자리에 허수부를 넣던 것을 NaN으로 \\
    옛 API & 2026-10-15 사용자 ``추천대로하자'': 폐기 예정 옛 API(\texttt{LSWTSolver.solve}, \texttt{EnergyFunction}, \texttt{SpinOptimizer})가 모두 거치는 \texttt{Diagonalizer.diagonalize\_w\_reg}에서, 원래 $H(k)$에 음의 고윳값이 있으면 모든 고윳값과 변환 행렬을 NaN으로 낸다(MAGSWT, \texttt{k-dependent} 모두). 영모드만 있는 상태는 그대로 \\
    검산 & \texttt{tests/test\_methods/test\_undefined\_values.py}(정의되는 경우와 안 되는 경우를 짝으로). 옛 해법과의 비교 시험은 임의 각도 상태에서 같은 $H(k)$인지를 보고, 에너지는 내지 않는지 본다 \\
    \bottomrule
  \end{tabularx}
  \slidesource{2026-10-15 결정}{contract}
\endgroup

\subsection{결정 기록 (D62): 최소 증강 스핀파(MAGSWT)를 켤 때만 쓴다}\label{decisions-d62-magswt}
\begingroup% frame-local settings, as in Beamer
  \framesubtitle{2026-10-20 사용자 확인(``그렇게하자''), 제안은 Claude. 잠정 승인.}
  \small\renewcommand{\arraystretch}{1.12}
  \begin{tabularx}{\textwidth}{@{}P{3.0cm}Y@{}}
    \toprule
    항목 & 결정 \\
    \midrule
    목적 & 고전적으로 불안정한 정류 상태의 $O(S)$ 에너지. 상이 고전 안정 영역 밖으로 넓어지는지(양자 escapism) 보려면 필요하다. D61은 그대로: 기본값에서 불안정하면 \texttt{LSWTError} \\
    API & \texttt{LSWTSettings(augmentation=None|"minimal"|}$\mu$\texttt{)}. \texttt{regularization="none"}과 $T=0$에서만. \texttt{compare\_states(..., augmentation=, refine="stationary")}는 상태 \texttt{"augmented"}와 $\mu$를 적고 안정 후보와 함께 순위를 매긴다 \\
    물리 & $\delta H=\mu\sum_i(S_i-\mathbf S_i\cdot\mathbf n_i)=\mu\sum_i a_i^\dagger a_i$(Holstein--Primakoff에서 정확). 고전 에너지는 그대로이고 $H(k)$ 대각에 $\mu$가 더해진다. 영점 에너지는 기존 식 $\tfrac12\sum(\varepsilon-\operatorname{tr}A)$을 증강된 행렬로 쓴다 \\
    $\mu$ 기준 & $\mu_{\min}=\max(0,-\min_k\lambda_{\min}H(k))$, 연속 $k$에서(D51과 같은 Nelder--Mead 탐색). $H+\mu\ge0$일 때만 진공이 증강된 이차 해밀토니안의 바닥 상태다. 논문의 기준($gH$ 고윳값이 실수, $\varepsilon^2>0$)은 음에너지 마그논을 허용하므로 비교값 \texttt{mu\_real\_spectrum}으로만 적는다. 옛 구현(\texttt{archive/legacy/modules/Tools/diagonalization.py} \texttt{prepare\_magswt})도 같은 $H\ge0$ 기준이었고, 다른 점은 메쉬 위 최솟값, 안정 상태에도 $\mu\ge10^{-9}$, 정류점·온도 확인 없음 \\
    영 모드 & $\mu=\mu_{\min}$에서 영 모드가 생긴다. $\mu_{\min}+10^{-9}\times$스펙트럼 크기를 쓰고, 그 에너지 차의 상한 $(\mu-\mu_{\min})\langle n\rangle$을 진단에 적는다($dE/d\mu=\langle n\rangle$) \\
    결과 & 진단 \texttt{energy\_bound="upper bound (MAGSWT, mu=...)"}. 의미: $\delta H\ge0$이므로 $E_{gs}(H)\le E_{gs}(H+\delta H)$가 엄밀하고, 숫자는 우변의 $O(S)$ 추정이다(엄밀한 변분 상한은 아님). 보손 수는 NaN, 띠·고유벡터를 쓰는 관측량 전부(띠, DOS, 구조인자, 중성자, 열, 위상, 두 마그논, 영모드 탐색)는 \texttt{require\_unaugmented}로 거부 \\
    한계 & 정류점이 아닌 상태는 지금처럼 NaN(증강해도 선형항은 남는다). 안정한 상태는 $\mu=0$이라 결과가 비트 단위로 같다 \\
    비교 대상 & Chernyshev arXiv:2509.08877 식 (25)--(27): 벌집 $J_1$--$J_3$ XXZ의 FM·ZZ $\mu$. 2026-10-20 확인: $J_3=0.30,0.35,0.60$에서 $10^{-7}$ 안으로 일치, XXZ 이방성과 무관, \texttt{mu\_real\_spectrum}도 같음(\texttt{tests/test\_methods/test\_magswt.py}). 다른 모형도 같은 날 확인: 같은 논문 식 (46)--(48)의 벌집 $J_1$--$J_2$ XXZ Néel·Iz·collinear $\mu$가 $J_2=0.1$--$0.5$, $\Delta=0,0.5,0.95$에서 $10^{-8}$ 안으로 일치. Coletta--Zhitomirsky--Mila arXiv:1212.3086 식 (6)--(7)의 정사각 $J_1$--$J_2$ 자기장 속 기운 Néel·줄무늬 $\mu$($-\min_k(A_k-|B_k|)$)가 일치하고, uuud($J_2=J_1/2$, $H=H_{\rm sat}/2$)와 등방 삼각 uud($H=H_{\rm sat}/3$)는 고전 점에서만 $\mu=0$. 다음: 같은 논문 그림 5·6 에너지와 DMRG(arXiv:2304.06062) 경계 \\
    \bottomrule
  \end{tabularx}
  \slidesource{2026-10-20 결정}{contract}
\endgroup

\subsection{결정 기록 (D63): 단위 세포를 고정한 정류점 탐색}\label{decisions-d63-stationary}
\begingroup% frame-local settings, as in Beamer
  \framesubtitle{2026-10-20 사용자 확인(``그렇게하자''), 제안은 Claude. 잠정 승인.}
  \small\renewcommand{\arraystretch}{1.12}
  \begin{tabularx}{\textwidth}{@{}P{3.0cm}Y@{}}
    \toprule
    항목 & 결정 \\
    \midrule
    목적 & MAGSWT는 정류점(안장점 포함)에만 쓸 수 있다(arXiv:2509.08877 4--5쪽). 결합 의존 항($J^{(3)}_{z\pm}$)이 있으면 대칭 무늬(UUD, 이중 지그재그)가 정류점이 아니다. \texttt{refine\_classical}은 에너지를 낮추므로 안장점에서 미끄러진다 \\
    API & \texttt{find\_stationary\_state(model, seed, conditions, torque\_tolerance=1e-10, max\_steps=100)} $\to$ \texttt{StationaryResult}(\texttt{state}, \texttt{converged}, \texttt{max\_torque}, \texttt{energy}, \texttt{saddle\_index}, \texttt{flat\_directions}, \texttt{max\_rotation\_from\_seed}, \texttt{steps}, \texttt{history}) \\
    방법 & 접선 좌표에서 $g=0$을 Newton 법으로 푼다(해석적 Hessian, 평평한 방향은 의사역행렬로 제외, 토크가 줄 때까지 걸음 반감, 실패하면 $-Hg$ 방향). 걸음마다 좌표를 다시 잡는다. $g$와 $H$가 모형 대칭에 공변이라 씨앗의 대칭이 유지된다 \\
    정의 안 됨 & 수렴하지 않으면 \texttt{converged=False}, \texttt{saddle\_index=-1}. 마지막 상태는 확인용일 뿐 LSWT 기준 상태가 아니다. \texttt{saddle\_index}는 세포 위의 Hessian 음의 방향 수이고, 연속 $k$ 안정성 판정은 D51 \\
    확인 & $J_3=0.45$, $J^{(3)}_{z\pm}=0.6$: UUD·이중 지그재그 씨앗(토크 0.3)에서 4걸음에 토크 $<10^{-10}$, 면 밖으로 기울어짐. $J^{(3)}_{z\pm}=0$의 UUD 안장점(지수 1)을 흔든 씨앗에서 같은 안장점으로 돌아옴 \\
    모형 & \texttt{models.bcao\_minimal}: Maksimov arXiv:2308.10672 식 (1). 결합 위상 벡터 $(\cos\phi_\alpha,\sin\phi_\alpha)$는 결합에 수직(입방 축 $\to$ 결정 축 사상에서 유도, 그림 2(a)는 읽지 못함), $\phi=0$은 $-y$ 방향 결합, $x$(자기장 방향)는 결합에 수직. 0장에서는 이 선택이 면내 전체 회전과 같아 에너지에 영향이 없다 \\
    비교 대상 & Maksimov 그림 1(a)·2(c)·3(a), 0장 고전 순서 ZZ $\to$ UUD $\to$ 이중 지그재그. 2026-10-20 고전 탐색(초격자 1, 2, 3, 4, 6줄, 줄 방향 하나): $J_3=0.34$에서 UUD는 $J^{(3)}_{z\pm}=0.1$, 이중 지그재그는 0.2부터로, 논문의 맞춤 점($J_3=0.34$, $J^{(3)}_{z\pm}=0.2$: 이중 지그재그, UUD 경계 가까이)과 맞는다. $J_3=0.45$에서는 UUD 구간이 약 0.42--0.55인데 본문은 UUD--이중 지그재그 전이를 ``0.9 부근''이라 적어 위치가 다르다. 그림을 읽지 못해 원인은 열린 항목 \\
    \bottomrule
  \end{tabularx}
  \slidesource{2026-10-20 결정}{contract}
\endgroup

\subsection{결정 기록 (D64): 발산은 부호 있는 무한대, 정의 안 됨은 NaN, 영 모드 보손 수는 탐지해 알린다}\label{decisions-d64-divergent-undefined}
\begingroup% frame-local settings, as in Beamer
  \framesubtitle{2026-10-20--21 사용자 지시(``보손수에 이상이 있는 것 같은 것이 자동적으로 탐지가 되었을 때, 사용자에게 알리고 진행하게 하자'', ``물리적으로 무한대로 발산하는 경우, 부호를 고려하여 무한대로 두고, 물리적으로 정의가 안되는 경우 nan으로 두자'', ``그렇게하자''), 적용은 Claude. D61을 고친다. 잠정 승인.}
  \small\renewcommand{\arraystretch}{1.12}
  \begin{tabularx}{\textwidth}{@{}P{3.0cm}Y@{}}
    \toprule
    항목 & 결정 \\
    \midrule
    원칙 & 극한이 있고 부호가 정해진 발산은 $\pm\infty$. 극한이 없거나(방향·메시·이동량 의존) 전개 밖이면 NaN과 경고. 발산하는 보손 수에서 얻는 모멘트 $S-\langle n\rangle$은 스핀파 전개 밖이라 NaN \\
    영 모드 보손 수 & 이상 탄성 항이 있는 골드스톤 점($B\ne0$)에서 $\langle n_k\rangle\sim\mu^{-1/2}\to+\infty$. 2차원 $T=0$ 적분은 유한하지만, 그 점을 포함한 메시 합은 이동량 $\mu$에 따라 달라지는 인공물이다(감사 CORE-01: $16^2,32^2,64^2$ MAGSWT 메시에서 $\langle n\rangle=61.9,15.6,4.05$). 정규화가 영 모드를 들어 올린 $k$에서 $H+\mu$와 $H+4\mu$의 점유를 비교해 달라지면 보손 수·모멘트를 NaN으로, 진단 \texttt{divergent\_zero\_mode\_occupation}, 선택지(Γ를 피하는 메시와 수렴 확인; 유한 토러스에서 대칭 골드스톤 모드를 빼고 크기 외삽; 모멘트 방향 작은 장에서 $L\to\infty$ 뒤 $h\to0$)를 경고에 적는다. 강자성($B=0$)은 $n=0$이 극한이라 그대로. 일반적인 영 모드 제외·외삽 기능은 만들지 않는다(사용자 결정) \\
    유한 온도 & 간격 없는 2차원 스펙트럼의 $t>0$ 보손 수는 $+\infty$(Mermin--Wagner), 모멘트·자화는 NaN, $F,U,S,C$는 유한. 정규화된 영 모드가 $k$ 집합에 있으면 $t>0$의 $F,S$는 $\ln\mu$에 의존해 NaN, $U,C$는 유지($\omega n_B\to t$) \\
    조화 전개 없음 & \texttt{thermal\_quantities}가 \texttt{stationary=False} 또는 $k\to-k$에 닫히지 않은 $k$ 집합에서 모든 열 양을 NaN으로(이전에는 $U,S,C,n,M$을 냈다, 감사 OBS-01). 구조인자, 스핀 상관, 결합 상관, 나선 구조인자, 두 마그논, 옷 입은 구조인자, 열 홀은 비정류 결과를 \texttt{require\_stationary}로 거부, Berry·Chern은 이차 부분 띠라는 경고만(OBS-02) \\
    나선 자화 & 회전 좌표 결과의 균일 자화는 실험실 좌표로: $Q$가 역격자 벡터가 아니면 회전축 성분만 남는다(OBS-03). \texttt{solve\_spiral\_lswt}가 저장하는 열 결과도 같은 규칙 \\
    구조인자 & 정규화가 들어 올린 영 모드 $q$도 D52 극한(Bragg)으로 다룬다. 이전에는 임의 이동량에서의 가중치를 냈다 \\
    비대각 성분 & 2026-10-21 사용자 선택(``part 인자''): \texttt{StructureFactor.spectrum}과 두 마그논 \texttt{spectrum}에 \texttt{part="real"|"imag"|"complex"}. 기본은 이전처럼 실수부. $S^{ab}$의 허수부는 반대칭(카이랄) 상관으로, 비편극 중성자 세기에는 안 보이지만 이전에는 버려졌다(감사 ARCH-01: $z$ 강자성의 $S^{xy}=\pm iS^{xx}$가 0). 대각합과 중성자 세기는 실수라 \texttt{"real"}만 받는다 \\
    비교 대상 & 정사각 Néel $S=1/2$: Γ를 피한 메시의 $\langle n\rangle\to0.1966$(Takahashi PRB 40, 2494 (1989); $64^2$에서 0.1926), Mermin--Wagner PRL 17, 1133 (1966). 원뿔 나선 자화 $=$ 초격자 자화(\texttt{tests/test\_methods/test\_zero\_mode\_occupation.py}) \\
    \bottomrule
  \end{tabularx}
  \slidesource{2026-10-21 결정}{contract}
\endgroup

\subsection{결정 기록 (D65): 새 API는 0.3에서 지울 옛 코드 없이 돈다}\label{decisions-d65-new-api-independent}
\begingroup% frame-local settings, as in Beamer
  \framesubtitle{2026-10-21 감사 2단계(사용자 ``다음 단계로 가자''), 적용은 Claude. 수치 결과는 그대로(전체 테스트 통과). 잠정 승인.}
  \small\renewcommand{\arraystretch}{1.12}
  \begin{tabularx}{\textwidth}{@{}P{3.0cm}Y@{}}
    \toprule
    항목 & 결정 \\
    \midrule
    목적 & D43이 0.3에서 지우는 \texttt{SpinSystem}, \texttt{LSWTSolver}, \texttt{EnergyFunction}, \texttt{SpinOptimizer}와 그것들만 쓰는 \texttt{LSWTHamiltonian}, 켈빈 \texttt{Thermodynamics}를 새 API가 부르지 않는다(감사 ARCH-03) \\
    상태 선택 & \texttt{LSWTZeroPointEnergy}/\texttt{LSWTHarmonicFreeEnergy}(constrained)가 \texttt{QuadraticBoseHamiltonian}과 자기 세포의 \texttt{BrillouinZone} 메시를 쓴다. 부드러운 모드 벡터는 같은 국소 좌표($e_\theta,e_\phi,s$)에서 만든다. \texttt{method="legacy"}는 \texttt{EnergyFunction}에 기대므로 폐기 경고 \\
    열역학 & \texttt{thermal\_quantities}와 \texttt{magnon\_free\_energy}가 무차원 함수(\texttt{bose\_occupation}, \texttt{bose\_free\_energy}, \texttt{bose\_entropy}, \texttt{bose\_heat\_capacity})로 모드 합을 직접 한다. 가짜 켈빈 $t/k_B$ 변환이 없어졌다(ARCH-14). $U=E_{cl}+E_{zp}+\sum w_k\varepsilon n_B/N_s$, $S=\sum w_k[(1+n)\ln(1+n)-n\ln n]/N_s$, $C=\sum w_k x^2/4\sinh^2(x/2)/N_s$ \\
    스핀 그림 & \texttt{plot\_spin\_configuration}이 모형·상태에서 자기 세포를 직접 만든다. onsite가 있는 모형도 그린다(ARCH-04). 결합선이 $r_i+(r_i-r_j)$로 그려지던 것을 $r_j$로 고쳤다(벌집처럼 반전 중심이 자리에 없는 격자에서 선이 빈 곳을 가리켰다) \\
    항 종류 & 새 함수 \texttt{require\_kinds}: \texttt{QuadraticBoseHamiltonian}, 비선형 전개, 나선 에너지, LT, 대칭 검사, 영 모드 대칭, \texttt{rotation\_symmetry}가 모르는 항 종류를 \texttt{NotImplementedError}로 거부한다. 이전에는 일부가 조용히 빼서 틀린 $H(k)$가 나올 수 있었다(ARCH-05) \\
    나선 onsite & 나선 에너지와 기울기가 onsite 항을 빠뜨렸다. 이제 $\kappa S^2 d^T\bar A d+(S/2)\operatorname{tr}A$, $\bar A$는 축 둘레 회전 평균(축 대칭이면 $\bar A=A$). $q=1/3$ 원뿔에서 초격자 고전 에너지와 $10^{-13}$ 안으로 일치 \\
    대칭 검사 & \texttt{CrystalSymmetry.check\_model}이 onsite 항을 본다: 상 $i$의 자리에서 $A'=R_sAR_s^T$(ARCH-02). C3를 깨는 $A=\mathrm{diag}(0.5,0,0)$은 위반, $D(S^z)^2$는 통과 \\
    검증 & \texttt{tests/test\_methods/test\_new\_api\_independence.py}: 지울 모듈을 쓸 수 없게 막은 하위 프로세스에서 \texttt{solve\_lswt}, \texttt{thermal\_quantities}, 상태 선택, 스핀 그림이 돈다. 이 셋은 사용자 선택으로 0.3 삭제 목록에 넣었다(D43 추가 행) \\
    \bottomrule
  \end{tabularx}
  \slidesource{2026-10-21 결정}{contract}
\endgroup

\subsection{결정 기록 (D66): 결과는 읽기 전용이고, 표준 JSON으로 저장되며, 계산 환경을 적는다}\label{decisions-d66-result-contract}
\begingroup% frame-local settings, as in Beamer
  \framesubtitle{2026-10-21 감사 3단계(사용자 ``이어서 진행하자.''), 적용은 Claude. 수치 결과는 그대로(전체 테스트 통과). 잠정 승인.}
  \small\renewcommand{\arraystretch}{1.12}
  \begin{tabularx}{\textwidth}{@{}P{3.0cm}Y@{}}
    \toprule
    항목 & 결정 \\
    \midrule
    목적 & 관측량이 확인하는 안전 표지(\texttt{extra["frame"]}, \texttt{extra["augmentation"]}, \texttt{header.diagnostics})를 계산 뒤에 바꿀 수 없게 하고(감사 ARCH-07), 저장한 결과가 계산을 다시 만들 정보를 갖게 한다(ARCH-08, C6) \\
    읽기 전용 & \texttt{LSWTResult}, \texttt{SpiralLSWTResult}, \texttt{ThermalResult}, \texttt{NLSWTResult}, 유사 골드스톤 결과, \texttt{EDResult}/\texttt{EDBlock}, \texttt{ClassicalSearchResult}, MC \texttt{ThermalAverages}, \texttt{ResultHeader}: 배열 필드는 읽기 전용 view, dict 필드는 \texttt{FrozenDict}(dict의 하위 클래스로 바꾸는 메서드가 \texttt{TypeError}). 넘겨 준 원래 배열은 그대로 쓸 수 있다. 바꾼 결과는 \texttt{dataclasses.replace}로 새로 만든다. 표지는 필드로 옮기지 않았다(공개 API 변경이라 토의 대상) \\
    비교 & 결과 클래스는 \texttt{eq=False}: \texttt{==}가 배열 때문에 오류를 내던 것을 동일성 비교로(ARCH-19) \\
    JSON & 표준 JSON(\texttt{allow\_nan=False}): NaN과 D64의 $\pm\infty$는 문자열 \texttt{"NaN"}, \texttt{"Infinity"}, \texttt{"-Infinity"}, 복소수는 \texttt{\{"real","imag"\}}. \texttt{load\_json}이 둘을 되돌린다(머리 형식 2판). LSWT JSON에 열 결과, 격자와 자기 격자, 위치, 국소 좌표계, 표지를 더했다. \texttt{NLSWTResult}, 유사 골드스톤 결과, \texttt{ClassicalSearchResult}에 \texttt{to\_json\_dict}. 결과 객체 자체를 JSON에서 되살리지는 않는다($H(k)$ 함수가 없으므로) \\
    계산 환경 & 머리에 python, numpy, scipy 버전. git 조회는 프로세스마다 한 번(결과마다 두 번 하던 9--12 ms, ARCH-09). git 저장소가 아닌 설치본은 pip이 남긴 커밋(\texttt{direct\_url.json}) \\
    seed & MC \texttt{ThermalAverages.settings}에 표본기의 seed와 실행 설정(seed가 numpy Generator이면 \texttt{"generator"}, 상태는 적지 않음). 고전 탐색은 이미 seed를 적는다 \\
    남은 것 & \texttt{StabilityReport}에 머리 없음, 검증 스크립트 출력에 머리 넣기, Langevin 표본의 seed는 적분기 객체에만 있음 \\
    검증 & \texttt{tests/test\_methods/test\_result\_contract.py}: 쓰기 시도가 막히고, NaN·$+\infty$·복소수가 엄격한 JSON 파서로 왕복하고, git을 다시 부르지 않고, MC seed가 남는다 \\
    \bottomrule
  \end{tabularx}
  \slidesource{2026-10-21 결정}{contract}
\endgroup

\subsection{결정 기록 (D67): $H(k)$는 결합마다 0이 아닌 원소만 저장하고, 대각화는 BLAS 1스레드를 권한다}\label{decisions-d67-sparse-assembly}
\begingroup% frame-local settings, as in Beamer
  \framesubtitle{2026-10-21 감사 4단계(사용자 ``다음으로 가자''), 적용은 Claude. 수치 결과는 비트 단위로 같음. 잠정 승인.}
  \small\renewcommand{\arraystretch}{1.12}
  \begin{tabularx}{\textwidth}{@{}P{3.0cm}Y@{}}
    \toprule
    항목 & 결정 \\
    \midrule
    목적 & \texttt{QuadraticBoseHamiltonian}이 결합마다 $(2N_s)^2$ 밀집 행렬 $P_b,Q_b$를 두어 메모리가 $N_s^3$으로 자라고, $H(k)$ 조립이 대각화보다 오래 걸렸다(감사 CORE-04) \\
    저장 & 결합마다 $P_b$, $Q_b$의 0이 아닌 원소 4개씩을 하나의 희소 행렬(행: $H$의 평탄화한 원소, 열: $2b$에 $P_b$, $2b+1$에 $Q_b$)에 둔다. 상수 $H_0$와 onsite는 밀집 그대로. $H(k)$와 해석적 1·2차 도함수는 희소 행렬과 위상 $(m,2n_b)$의 곱 하나 \\
    대각화 & 밀집 Colpa 그대로($N_s^3$). 셀이 작아 희소 고유값 풀이의 이득이 없다 \\
    측정 & 정사각 FM 초격자, 1스레드, 16 k점의 $H$와 도함수: $N_s=144$에서 조립 메모리 1.53\,GB → 2\,MB, 시간 16.7\,s → 0.34\,s. $N_s=400$, $2\times2$ 메시 \texttt{solve\_lswt} 5.4\,s, 최대 메모리 273\,MB(이전 방식은 결합 행렬만 약 16\,GB로 추산). 이전 구현과 $H$, 도함수 차이 0(삼각 격자 무작위 비공선 상태 포함) \\
    스레드 & 작은 행렬을 k마다 푸는 구조라 다중 스레드 BLAS가 경합만 한다(감사 CORE-05: 1스레드 0.02\,s 대 경합 4스레드 9\,s). README ``Performance''에 \texttt{OMP\_NUM\_THREADS=1}과 프로세스 병렬을 권하고, CI 테스트도 1스레드로 돌린다. \texttt{threadpoolctl}은 의존성으로 넣지 않았다 \\
    \bottomrule
  \end{tabularx}
  \slidesource{2026-10-21 결정}{contract}
\endgroup

\subsection{결정 기록 (D68): 정지성은 상대 토크로 판정하고, 안정성은 무엇을 검사했는지 적고, 검증 기준을 정한다}\label{decisions-d68-relative-stationarity}
\begingroup% frame-local settings, as in Beamer
  \framesubtitle{2026-10-21 감사 질문 3, 5, 6(사용자 ``다음은 추천대로 진행해''), 적용은 Claude. 잠정 승인.}
  \small\renewcommand{\arraystretch}{1.12}
  \begin{tabularx}{\textwidth}{@{}P{3.0cm}Y@{}}
    \toprule
    항목 & 결정 \\
    \midrule
    정지성 & 허용치 \texttt{stationarity\_tolerance}(기본 $10^{-8}$)는 상대 토크 $\max_i|S_i\times h_i|/\max_i S_i|h_i|$에 적용한다(감사 SPEC-13). 스핀과 국소장 사이 가장 큰 어긋남 각의 사인이고, 결합·장·$S$를 함께 늘려도 그대로다. 이전 절대값 $10^{-8}\,E_0$는 $J$가 작은 단위에서 어긋난 상태를 정지 상태로 받았다($J=10^{-9}$, 어긋남 $10^{-4}$ rad). 머리에 \texttt{max\_torque}($E_0$)와 \texttt{relative\_torque}를 함께 적는다. \texttt{solve\_lswt}, \texttt{local\_stability}, \texttt{solve\_spiral\_lswt}에 적용. 나선의 피치 기울기 $|dE/dq|$는 같은 장 크기에 격자 길이 $\sqrt{|\det a|}$를 곱한 값으로 나눈다 \\
    안정성 표시 & \texttt{solve\_lswt}는 연속 $k$ 안정성을 기본으로 검사하지 않고, 머리 \texttt{diagnostics["stability\_checked"]}에 ``computed momenta only; local\_stability checks continuous k''(MAGSWT면 ``continuous k'')를 적는다(SPEC-10). 연속 검사는 \texttt{local\_stability}(D51) \\
    검증 기준 & AGENTS.md 작업 원칙 3: 구현은 기준값 대비 검증 뒤 다음 단계로. 기준값은 닫힌 식, 작은 ED, Toth--Lake 관례의 독립 계산, magnopy 같은 다른 패키지. 옛 코드 수치는 공선 상태에서만(옛 코드는 비공선 상태에서 $B$와 $B^\dagger$를 바꿈, \texttt{hamiltonian\_convention.tex})(SPEC-17) \\
    검증 & \texttt{tests/test\_methods/test\_relative\_stationarity.py}: 상대 토크가 $J$, $S$ 배율과 무관하고, 판정이 에너지 단위에 의존하지 않는다 \\
    \bottomrule
  \end{tabularx}
  \slidesource{2026-10-21 결정}{contract}
\endgroup

\subsection{결정 기록 (D69): 메시가 풀지 못하는 1/S 상쇄는 NaN으로 낸다}\label{decisions-d69-goldstone-cancellation}
\begingroup% frame-local settings, as in Beamer
  \framesubtitle{2026-10-21 이론적 신뢰성 확인(사용자 ``신뢰성 확인하자''), 적용은 Claude. 잠정 승인. 비교 대상: Chernyshev--Zhitomirsky PRB 79, 144416 (2009) 식 (56).}
  \small\renewcommand{\arraystretch}{1.12}
  \begin{tabularx}{\textwidth}{@{}P{3.0cm}Y@{}}
    \toprule
    항목 & 결정 \\
    \midrule
    현상 & 비공선 상태에서 자기 영역 중심으로 접힌 Goldstone 모드(120° 상태의 $\pm Q$ 모드) 근처에서 Hartree--Fock+tadpole 항과 $\mathrm{Re}\,\Sigma_3$가 각각 $1/|k|$로 커지고 합은 $|k|$로 작아진다. 둘 다 메시 합이라 $1/N$ 크기의 어긋남이 비 $\rho=\max(|\delta E_{\mathrm{HF}}|,|\mathrm{Re}\,\Sigma_3|)/\omega$배로 커진다. 삼각 $S=1/2$, $N=24$에서 격자 밖 $k$의 비 $\bar\varepsilon/\varepsilon$이 0.4--3.3(기대 1.08)으로 경고 없이 나왔다 \\
    격자 위 $k$ & $k=(j_1 b_1/N_1, j_2 b_2/N_2)$이면 3마그논 합의 $k-q$가 모두 메시($-1/3$ 칸 이동)에 놓여 상쇄가 정확하다. 이 경우 식 (56)의 속도 비 $1-0.11488/2S$, $1+0.08285/2S$와 2\,\% 안에서 맞는다 \\
    규칙 & \texttt{magnon\_energies}(격자 밖 $k$)와 \texttt{spectral\_function}(모든 $k$)에서 $\rho>N_{\min}/4$인 띠의 에너지(와 스펙트럼)는 NaN, 경고에 선택지(격자 위 $k$, 더 촘촘한 메시, 영모드에서 먼 $k$)를 적는다. 유한 토러스는 정확한 섭동론이라 제외. $N/4$는 $N=24, 48$의 측정(\texttt{data-space/verification/261021-theory-reliability})에서 정한 경험 기준이다 \\
    남은 일 & 상쇄를 해석적으로 빼는 방법(회전 틀 또는 Goldstone 발산 분리)은 하지 않았다. 띠가 두 마그논 연속체 안에 있을 때 $\eta=0$의 메시 주값도 메시에 민감하다(삼각 $N=24$와 48에서 0.87과 1.18). 이것은 사용자 결정으로 남긴다 \\
    검증 & \texttt{tests/test\_methods/test\_reliability\_benchmarks.py} \\
    \bottomrule
  \end{tabularx}
  \slidesource{2026-10-21 결정}{contract}
\endgroup

\subsection{결정 기록 (D70): 고전 탐색은 LT 하한을 기준으로 적응적으로 다시 돈다}\label{decisions-d70-adaptive-restarts}
\begingroup% frame-local settings, as in Beamer
  \framesubtitle{2026-10-21 이론적 신뢰성 확인, 사용자 승인(``그렇게하자''). 적용은 Claude. 잠정 승인. 비교 대상: 같은 초격자의 LT 하한, BCAO 시드별 에너지(\texttt{data-space/verification/261021-theory-reliability}).}
  \small\renewcommand{\arraystretch}{1.12}
  \begin{tabularx}{\textwidth}{@{}P{3.0cm}Y@{}}
    \toprule
    항목 & 결정 \\
    \midrule
    현상 & 차분 진화 한 번이 국소적으로 안정한 국소 최소에 갇힌다. BCAO $J_3=0.45$, $J_{z\pm}=0.5$, 셀 $[[1,0],[0,6]]$에서 기본 시드는 UUD($-0.303614$, \texttt{local\_stability} 판정 stable), 다른 시드는 $-0.304193$. 그래서 국소 안정성만으로는 다시 돌릴지를 정할 수 없다 \\
    하한 & 초격자에 맞는 운동량 $q$($q\cdot R\in 2\pi\mathbb{Z}$)에서 $J(q)$의 최저 고유값 + 온사이트 상수는 그 초격자 위 모든 상태 에너지의 엄밀한 하한이다(약한 제약 LT). 도달하면 전역 최소가 증명된다. 0이 아닌 장에서는 정의하지 않는다(\texttt{None}) \\
    규칙 & \texttt{classical\_search(restarts="adaptive")}가 기본. 시드 \texttt{seed}, \texttt{seed+1}, … 순서로 돌려 (1) 하한에 도달하면, (2) 연속 2번 최저값이 $10^{-8}$보다 더 내려가지 않으면, (3) 8번이면 멈춘다. 정수를 주면 그 수만큼 돈다(\texttt{restarts=1}이 옛 동작). 결과 \texttt{settings}에 시드별 에너지, 폭, 하한, 도달 여부, 멈춘 이유를 적는다 \\
    비용 & 하한에 닿는 상(FM, Néel, 120°, 지그재그 등)은 한 번. 결합 의존 항이나 장이 있으면 보통 세 번 이상 \\
    남은 일 & 하한에 닿지 않으면 전역 최소의 증명은 없다. 시드 폭은 증거일 뿐이다. 다중 $q$ 상태와 장 속에서는 다른 하한(일반화 LT)이 없다 \\
    검증 & \texttt{tests/test\_methods/test\_classical\_search.py} \\
    \bottomrule
  \end{tabularx}
  \slidesource{2026-10-21 결정}{contract}
\endgroup

\subsection{결정 기록 (D71): 연속체 안의 $\eta=0$ 메시 에너지는 NaN으로 낸다}\label{decisions-d71-eta-zero-continuum}
\begingroup% frame-local settings, as in Beamer
  \framesubtitle{2026-10-21 이론적 신뢰성 확인, 사용자 승인(``η=0 주값 처리 … 진행하자'', 추천안 NaN). 적용은 Claude. 잠정 승인. 비교 대상: \texttt{spectral\_function}의 주값 에너지(삼각형 적분 밀도의 Hilbert 변환).}
  \small\renewcommand{\arraystretch}{1.12}
  \begin{tabularx}{\textwidth}{@{}P{3.0cm}Y@{}}
    \toprule
    항목 & 결정 \\
    \midrule
    현상 & 띠 $\omega_n(k)$가 두 마그논 붕괴 연속체 안에 있으면 $\eta=0$의 메시 합 $\Sigma_3(k,\omega_n)$이 극을 지나며, 실수부가 주값이 아니라 메시마다 다른 수가 된다. 삼각 $S=1/2$의 같은 $k$에서 $N=24$, 48이 0.87과 1.18 \\
    규칙 & 열역학 극한 기하에서 \texttt{magnon\_energies(eta=0)}의 \texttt{total}은 그 띠에서 NaN, 경고에 선택지(\texttt{spectral\_function}의 \texttt{energy}, 양의 $\eta$, 유한 토러스)를 적는다. 판정은 메시의 짝 $(q,k-q)$ 중 붕괴 꼭짓점이 0이 아닌 것($|D|^2>10^{-10}\max$)이 $\omega_{n_1}(q)+\omega_{n_2}(k-q)<\omega_n(k)$를 만족하는가이다. 숨은 대칭으로 꼭짓점이 0인 통로는 세지 않는다 \\
    확인 & 삼각 $S=1/2$, $k$ = Γ–K 중간: 예전 메시 값 1.0885는 주값 1.1778(메시 24, 48과 격자 96, 192에서 $10^{-3}$ 안)과 8\,\% 달랐다. $k\to0$ Goldstone 모드도 $\pm Q$ 둘로 붕괴할 수 있어 NaN이 되며, 그 속도 비는 주값으로 Chernyshev--Zhitomirsky 식 (56)과 2\,\% 안. 연속체 밖 띠는 메시 값이 낫다(Hilbert 변환은 \texttt{levels}에 따라 수 \% 흔들림: 0.7355 대 메시 0.7313) \\
    검증 & 연속체 밖 띠는 $N=24$, 48에서 $2\times10^{-3}$ 안에서 같다. \texttt{tests/test\_methods/test\_reliability\_benchmarks.py} \\
    \bottomrule
  \end{tabularx}
  \slidesource{2026-10-21 결정}{contract}
\endgroup

\subsection{결정 기록 (D72): 자기무모순 one-loop(SCOL)는 3차 꼭짓점이 없는 상태에서만}\label{decisions-d72-self-consistent-one-loop}
\begingroup% frame-local settings, as in Beamer
  \framesubtitle{2026-10-21 사용자 요청(``self-consistent를 일반 패키지로 더하자''). 범위와 규칙은 Claude 추천, 잠정 승인. 비교 대상: Keselman et al., arXiv:2512.02150 End Matter와 SM III절(SCOL), Chubukov--Golosov 1991(큰 $S$ 극한).}
  \small\renewcommand{\arraystretch}{1.12}
  \begin{tabularx}{\textwidth}{@{}P{3.0cm}Y@{}}
    \toprule
    항목 & 결정 \\
    \midrule
    방법 & \texttt{solve\_scswt}: $H_\mathrm{SC}(k)=H_2(k)+\delta H_\mathrm{HF}$, 두 점 함수 $\langle a^\dagger a\rangle$, $\langle aa\rangle$를 $H_\mathrm{SC}$의 진공에서 다시 구해 바뀌지 않을 때까지 되풀이(Keselman 식 (S53)--(S54)). 고정점은 Bogoliubov 진공들 위 가우스 에너지 $\langle H_0+H_2+H_4\rangle$의 정류점이고, 이 값을 바닥 에너지로 낸다. 마그논 에너지는 $H_\mathrm{SC}$의 Bogoliubov 에너지 \\
    범위 & 3차 꼭짓점이 없는 상태($H_3=0$: UUD 같은 공선 판, 대칭축을 따른 공선 상태)만. $H_3\neq0$이면 $\Sigma_3$와 tadpole이 HF와 같은 차수로 마그논 에너지를 옮기므로, HF만 자기무모순으로 다루는 것은 일관된 근사가 아니다. 이때 \texttt{LSWTError}와 선택지(\texttt{solve\_nlswt}) \\
    장 스캔 & \texttt{frozen=}: 다른 장의 결과에서 $\delta H_\mathrm{HF}$를 그대로 쓴다(4차 항은 장과 무관). $S^z_\mathrm{tot}$가 보존되면 장 항이 $H_\mathrm{SC}$와 교환해 진공이 그대로이므로, 고정은 근사가 아니라 정확하다(Keselman의 ``약하게 의존한다'' 가정이 필요 없음). \texttt{initial=}: 다른 장의 해에서 반복을 시작(LSWT 진공이 없는 곳) \\
    NaN 규칙 & 고정 모드에서 메시 위 $H_\mathrm{SC}$가 양의 정부호가 아니면(간극 닫힘 너머) 에너지와 보손 수는 NaN과 경고. \texttt{magnon\_energies}는 그 $k$에서만 NaN. 자기무모순 반복이 불안정한 $H$에 닿거나 수렴하지 않으면 오류 \\
    확인 & 정사각 Néel: HF가 $H(k)$를 균일하게 늘리므로 고정점 = one-loop($10^{-6}$). 큰 $S$ 정확 대각화(XYZ 편극상, $2\times2$ 토러스, $S=1$--8): $(E_\mathrm{ED}-E_\mathrm{SC})S$, 간극 차$\times S$가 유한하게 수렴. 정류성: SC 진공에서 에너지 기울기 $10^{-8}$, LSWT 진공에서 $10^{-2}$. UUD: $n_\downarrow=2n_\uparrow$ 정확히, $S\to\infty$에서 경계가 CG $0.50$, $1.29$로 \\
    결과 & 삼각 UUD $S=1/2$: $h_{c1}=1.2665J$, $h_{c2}=2.0517J$(one-loop 1.25, 2.15), $|\langle S^z\rangle|$ 위 0.4168, 아래 0.3336. Keselman의 SCOL·DMRG 값은 그림으로만 있어 숫자로 대조하지 못했다(미확인) \\
    검증 & \texttt{tests/test\_methods/test\_self\_consistent.py}, 표 \texttt{data-space/verification/261021-self-consistent/tables.txt} \\
    \bottomrule
  \end{tabularx}
  \slidesource{2026-10-21 결정}{contract}
\endgroup

\subsection{결정 기록 (D73): 3차 꼭짓점의 평균장 처리(Bader식 각도 재규격화)}\label{decisions-d73-mean-field-cubic}
\begingroup% frame-local settings, as in Beamer
  \framesubtitle{2026-10-21 사용자 요청(``bader 확장 식을 진행하자''; 앞서 결정 카드에서 ``Bader식은 다음 PR''). 방식과 규칙은 Claude 추천, 잠정 승인. 비교 대상: Bader et al., arXiv:2512.08768 SM S3절(식 S13--S15), \texttt{solve\_nlswt}의 tadpole과 3차 2차 섭동 에너지.}
  \small\renewcommand{\arraystretch}{1.12}
  \begin{tabularx}{\textwidth}{@{}P{3.0cm}Y@{}}
    \toprule
    항목 & 결정 \\
    \midrule
    방법 & \texttt{SCSWTSettings(cubic="mean\_field")}: 스핀 방향 $\theta$도 변수. 가우스 진공에서 $H_3$를 축약하면 선형 항이 남고, $H_1(\theta)$와 합친 힘 $H_1+\langle H_3\rangle_\mathrm{lin}$이 0이 될 때까지 방향을 돌린다(뉴턴 단계 $\langle\psi\rangle=-H_\mathrm{SC}(0)^{+}\zeta$, $\langle s^+\rangle=\sqrt{2S}\langle a\rangle$). 동시에 $H_\mathrm{SC}=H_2(\theta)+\delta H_\mathrm{HF}[H_4(\theta)]$를 자기무모순으로. 에너지 $\langle H_0+H_2+H_4\rangle(\theta)$. 모든 차수를 현재 $\theta$에서 전개(Bader는 $H_4$를 고전 각도에서; 차이는 더 높은 차수) \\
    빠지는 것 & 3차 꼭짓점의 2차 과정($E_\mathrm{cubic}$, 마그논의 $\Sigma_3$). 바닥 에너지에서 HF와 같은 $S^0$ 차수다. 가우스 상태는 세 마그논 상관을 담지 못한다. 자기무모순 진공에서 계산한 그 값을 \texttt{energies.three\_magnon}으로 따로 내고 \texttt{total}에는 넣지 않는다 \\
    골드스톤 & one-loop에서 골드스톤 정리는 HF + tadpole + $\Sigma_3(0,0)=0$으로 성립한다. $\Sigma_3$가 없으면 $k=0$ 영 모드가 가짜 간극을 얻는다(Hohenberg--Martin; 상대 크기 $S^{-1/2}$; 정사각 Heisenberg 기울어진 반강자성, $S=1$에서 간극 $1.6J$). 고전 $H(0)$에 영 모드가 있고 $H_3\neq0$이면 경고, \texttt{diagnostics["spurious\_gap"]}, \texttt{magnon\_energies}는 NaN. 에너지는 내되 $S^0$ 차수 오차가 $1/S$ 대신 $S^{-1/2}$ 상대 \\
    시작 상태 & 고전 정류 상태가 아니어도 된다($H_1$이 힘에 들어감; 그때 \texttt{one\_loop}는 NaN). \texttt{initial=}: 같은 모형·자기 단위포·운동량의 평균장 결과에서 방향과 진공을 이어받는다(장 스캔). \texttt{frozen=}은 안 됨(방향이 장에 따라 바뀜) \\
    확인 & $H_3=0$이면 \texttt{cubic="refuse"}와 같다($10^{-10}$). 간극 있는 기울어진 XYZ 반강자성: 각도 이동$\times S$가 \texttt{solve\_nlswt} 응축 값으로, $E_\mathrm{MF}-(E_\mathrm{1L}-E_\mathrm{cubic})$이 $1/S$로 줄어든다. 큰 $S$ ED(3$\times$1 BCAO 토러스): $E_\mathrm{ED}-E_\mathrm{MF}\to E_\mathrm{3m}$, $E_\mathrm{ED}-E_\mathrm{MF}-E_\mathrm{3m}$은 $1/S$로 줄어든다(B13). $S=1$에서는 평균장이 one-loop보다 ED에 가깝다 \\
    검증 & \texttt{tests/test\_methods/test\_self\_consistent.py}(평균장 4개), B13, 표 \texttt{data-space/verification/261021-self-consistent/tables.txt} 3절 \\
    \bottomrule
  \end{tabularx}
  \slidesource{2026-10-21 결정}{contract}
\endgroup

\subsection{결정 기록 (D74): 외부 물리 검토(2026-10-23)에서 확인한 오류 수정}\label{decisions-d74-review-fixes}
\begingroup% frame-local settings, as in Beamer
  \framesubtitle{2026-10-23 사용자 지시(``다시한번 확인해보고, 문제가 명확하면 수정해''). ChatGPT 검토 P1--P9의 지적 가운데 Claude가 다시 재현한 열 가지. 방식은 Claude 추천, 잠정 승인.}
  \small\renewcommand{\arraystretch}{1.12}
  \begin{tabularx}{\textwidth}{@{}P{3.0cm}Y@{}}
    \toprule
    항목 & 결정 \\
    \midrule
    교환 행렬 & \texttt{xxz}, \texttt{xxz\_with\_soc}는 실수 행렬을 만든다(정수 입력에서 $\Gamma$가 0으로 잘리던 오류) \\
    나선 에너지 & \texttt{spiral\_energy}: 결합·자리·Zeeman 항을 격자가 실현하는 위상 $2\pi q\cdot R$로 평균한다($q$ 정수: 한 점, $2q$ 정수: $0,\pi$, 그 밖: 등간격 8점이 2차 조화까지 정확). U(1) 모형에서는 그대로. 확인: 비U(1) 모형에서 초격자 \texttt{classical\_energy}와 $10^{-13}$ \\
    영모드 띠 & \texttt{band\_structure}의 준정부호 분기: $\eta H$ 고유값 가운데 큰 $N_s$개가 입자 에너지 $\omega_n(k)\ge0$, 나머지가 $-\omega_n(-k)$. 모듈러스를 둘씩 묶던 방식은 $\omega(k)\neq\omega(-k)$에서 틀린다 \\
    상태 비교 & \texttt{compare\_states}: LSWT가 비정류라고 판정한(D68) 후보는 \texttt{not stationary}, 순위에서 뺀다. 조화 에너지가 NaN인 후보는 \texttt{stable}로 두지 않는다 \\
    $n>S$ & \texttt{solve\_lswt}: $\langle n_i\rangle>S_i$이면 경고, \texttt{diagnostics["max\_boson\_number\_per\_spin"]}. 순서 모멘트는 이미 NaN. 에너지는 그대로 낸다(값을 NaN으로 할지는 사용자 결정 대기) \\
    축퇴 띠(1/S) & \texttt{NLSWTResult.magnon\_energies}: 정확히 축퇴한 띠의 보정은 $P(\delta H+\Sigma_3(\omega_0))P$의 고유값(축퇴 섭동론). \texttt{spectral\_function}: 축퇴 공간의 행렬 밀도(\texttt{cubic\_spectral\_density\_matrix}, $\rho_{nm}\propto D_n^*D_m$)로 같은 행렬을 만들어 대각화, Dyson 모드의 $\Sigma$도 그 고유 기저에서. 대칭으로 보호된 축퇴는 삼각형 격자 오차만큼($10^{-4}$) 갈라질 수 있다. 확인: $JS$가 같은 독립 정사각 AFM 두 층($S=1/2,1$)이 각 층의 값으로 갈라진다 \\
    종방향 연속체 & \texttt{two\_magnon\_continuum}($t=0$): 두 마그논이 함께 부드러워지는 $q$(Néel Bragg점)에서 가중치가 $|k|^{-2}$, 2차원 합이 로그 발산한다. 영모드 주위 원에서 지수 $p\ge1.9$이면 그 대각 성분의 \texttt{static}과 넓힌 \texttt{spectrum}은 $+\infty$(비대각 NaN), 경고. 확인: 메시를 두 배로 할 때마다 $\ln2/\pi=0.2206$씩 증가; $q=0$은 $S^z_\mathrm{tot}$ 보존으로 0 \\
    평탄한 쌍 에너지 & 삼각형 적분(\texttt{triangle\_density}): 쌍 에너지가 삼각형 위에서 상수이고 무게가 있으면 측도에 $\delta$ 원자가 있어 밀도로 나타낼 수 없다. NaN과 경고(유한 토러스를 쓰라고 안내). 황금률 적분은 그 상수가 0일 때만 NaN \\
    Chern 수 & \texttt{chern\_numbers}: Kubo--FHS가 일치해도 연속 $k$에서 띠 간격을 국소 최소화해, 간격이 닫히거나 메시 최솟값의 0.5배(\texttt{TOPOLOGY\_GAP\_RESOLUTION}) 아래이면 NaN과 경고. 반례: 벌집 강자성 Dirac점 합병($J_1=2$, $D=0$)과 $J_1=1.999$, $D=10^{-3}$(참값 $(+1,-1)$, 메시 12에서 $(0,0)$) \\
    BZ & \texttt{high\_symmetry\_points}, \texttt{zone\_boundary}: 역격자 기저를 먼저 Gauss 축약한다(\texttt{A=[[1,0],[8,1]]}에서 BZ가 26배였음) \\
    검증 & \texttt{tests/test\_methods/test\_physics\_review\_2026\_10\_23.py} \\
    \bottomrule
  \end{tabularx}
  \slidesource{2026-10-23 결정}{contract}
\endgroup

\subsection{결정 기록 (D75): 통제되지 않은 진공과 메시 사이 불안정}\label{decisions-d75-vacuum-stability}
\begingroup% frame-local settings, as in Beamer
  \framesubtitle{2026-10-24 사용자 지시(``그래 그렇게 진행해'', 급한 항목 1·2번). ChatGPT 검토 P10의 남은 P1 묶음 R05(진공 전파)와 안정성 오인증 가운데 Claude가 재현한 것. 방식은 Claude 추천, 잠정 승인.}
  \small\renewcommand{\arraystretch}{1.12}
  \begin{tabularx}{\textwidth}{@{}P{3.0cm}Y@{}}
    \toprule
    항목 & 결정 \\
    \midrule
    $n>S$ 원칙 & $\langle n_i\rangle(t)>S_i$이면 모멘트가 뒤집혀 그 상태 둘레의 스핀파 전개가 통제되지 않는다. 진공(또는 열 앙상블)으로 만든 양은 NaN과 경고(\texttt{vacuum\_beyond\_spin}). 띠(고전 미소 진동수)는 둔다. LSWT 바닥 에너지 값은 사용자 결정 대기(D74), \texttt{bond\_correlations}의 에너지도 같은 처리 \\
    적용 & \texttt{structure\_factor}, \texttt{spiral\_structure\_factor}(가중치·탄성), \texttt{spin\_correlation}, \texttt{bond\_correlations}(상관), \texttt{two\_magnon\_continuum}, \texttt{thermal\_hall}($t$별, \texttt{thermal\_quantities}와 같은 기준), \texttt{LSWTHarmonicFreeEnergy}, \texttt{solve\_nlswt}(HF·3차·tadpole 에너지와 모든 $1/S$ 마그논 양; HF 평균이 진공 기댓값을 직접 쓴다), \texttt{structure\_factor\_s0}, \texttt{structure\_factor\_states} \\
    MAGSWT 띠 & \texttt{LSWTResult.bands()}도 \texttt{require\_unaugmented}: $\mu$만큼 이동한 띠를 모형의 띠로 내지 않는다. JSON에서는 \texttt{null} \\
    메시 사이 음수 & \texttt{scan\_zero\_modes}: $\lambda_\mathrm{min}/\mathrm{scale}<-$\texttt{zero\_tolerance}는 영모드가 아니라 \texttt{unstable}(\texttt{has\_unstable}). 열역학량·구조 인자($t>0$)·연속체($t>0$)·열 홀·자화 곡선은 NaN과 경고. 반례: $J_1=1$, $J_2=0.7$ 사슬과 $y$ 강자성, $h=h_\mathrm{sat}-0.01$에서 정확한 최저 $\omega=-0.01$, 메시 8에서는 최저 0.067 \\
    SCSWT frozen & 열역학 극한에서 $H_\mathrm{SC}(k)$의 최저 고유값을 특수점($0,1/2,1/3,2/3$)과 메시 최저점에서 연속 최소화한다. 음수이면 frozen은 NaN과 경고, 자기무모순은 \texttt{LSWTError}. 반례: UUD $S=1/2$, $h_{c1}-10^{-4}$(1/3 이동 메시에 $k=0$이 없다) \\
    대칭 허용치 & \texttt{\_symmetry\_axes}: 허용치를 1 E0가 아니라 결합 크기에 상대로. 모든 결합을 $c$배 해도 판정이 같다. 반례: $S=1$ Néel spinodal, $c=10^{-12}$에서 \texttt{stable} \\
    4차 표본 & 2차원 이상 우연 영공간에서 표본 방향이 모두 오르면, 가장 작은 진폭에서 방향을 연속 최소화하고 그 방향도 판정한다. 반례: 4차 형식 $[x_1^4+4x_2^4-4.1x_1^2x_2^2]/2$, 8방향 모두 양수지만 $-(2/45)a^4$ 방향이 있다 \\
    LT 하한 & \texttt{lambda\_min}은 찾은 최솟값이며 전역 최소라는 인증이 없으면 하한이 아니다(\texttt{diagnostics["global\_minimum\_certified"]=False}, 설명 수정). 결합 길이가 메시의 절반 이상이면 경고 \\
    상태 선택 & \texttt{select\_on\_manifold}: 고른 상태를 \texttt{local\_stability}(연속 $k$)로 확인해 불안정이면 \texttt{not\_soft}. 반례: compass 모형 $j_3>a/4$에서 Tetra $N=4$ 메시가 놓친 $q^2$ 불안정 \\
    검증 & \texttt{tests/test\_methods/test\_physics\_review\_2026\_10\_24.py}; 고온 열 홀 기댓값 갱신(튜토리얼 5의 $T=1$은 NaN) \\
    \bottomrule
  \end{tabularx}
  \slidesource{2026-10-24 결정}{contract}
\endgroup

\subsection{물리 검토 원칙 (2026-10-15)}\label{decisions-physics-review-principles}
\begingroup% frame-local settings, as in Beamer
  \framesubtitle{2026-10-15 사용자 지시(``물리검토는 원칙을 세우고, 추천대로 진행해.'' ``실제 엄밀한 검토는 관련 계산을 할 때, 논문결과와 비교하면서 검토할거야''), 원칙은 Claude 제안.}
  \small\renewcommand{\arraystretch}{1.12}
  \begin{tabularx}{\textwidth}{@{}P{3.0cm}Y@{}}
    \toprule
    항목 & 원칙 \\
    \midrule
    잠정 승인 & D50--D60은 추천안대로 쓴다. 이것은 사용자의 엄밀한 물리·수학 검토가 아니다. 엄밀한 검토는 그 결정을 쓰는 계산을 할 때 아래 비교 대상의 논문 결과와 맞춰 보며 한다. 그때 어긋나면 결정을 다시 연다 \\
    판정 기준 & (1) 주장은 유도나 정확히 알려진 극한에 근거하고, 독립된 수치 검산(닫힌 식, 다른 방법, 큰 $S$ ED)이 하나 이상 있어야 잠정 승인한다. (2) 정의되지 않거나 비물리적인 값은 NaN(D61). (3) 수치 바닥값에 기댄 판정(``평평'', ``수렴'')은 증명이 아니라 수치 판정이라고 적는다. (4) $1/S$로 통제되지 않는 계산(Dyson 재합 등)은 기본값이 아니라 선택 사항과 진단값으로 둔다. (5) 관례로 정한 나눔(붕괴하는 선과 연속체)은 관례라고 밝히고 정확도 밖의 차이만 허용한다 \\
    D50, D51 & 수치 방법과 판정 기준. 비교 대상: 초격자 판정과 같은 결과(D51 시험), NBCP Y 고전 궤도 \\
    D52 & 비교 대상: 120° A4 닫힌 식, 정사각 Néel의 $q\to0$ 극한 $t/8J$(고전 가로 감수율) \\
    D54 & 비교 대상: 장 속 삼각 하이젠베르크의 고전 축퇴와 order by disorder(Chubukov--Golosov 1991), 쉬운 축 Néel의 1차 spin flop \\
    D55 & 비교 대상: $J_1$--$J_2$ 사슬 나선, 삼각 120° 초격자 판정 \\
    D56, D60 & 비교 대상: 정사각 Néel 종방향 연속체의 문헌 세기, 상세 균형 \\
    D57 & 비교 대상: 기울어진 정사각 AFM 붕괴 시작 $h^*\approx0.76h_c$(Zhitomirsky--Chernyshev 1999), 삼각 AFM 붕괴율(Chernyshev--Zhitomirsky 2009) \\
    D58, D59 & 비교 대상: 삼각 AFM의 $1/S$ 스펙트럼과 $S(q,\omega)$(Chernyshev--Zhitomirsky 2009; Mourigal 외 2013). 붕괴하는 띠의 무게는 선폭 평균값(관례)이라 그 비교에서 특히 본다 \\
    \bottomrule
  \end{tabularx}
\endgroup

\subsection{세부 규칙: 2026-09-29 결정}\label{detail-rules}
\begingroup% frame-local settings, as in Beamer
  \framesubtitle{D12 · 1차 범위를 규약으로 옮기며 추가한 규칙에 대해 결정한 내용이다.}
  \small
  \begin{tabularx}{\textwidth}{@{}p{0.4cm}YP{4.9cm}@{}}
    \toprule
    & 결정 & 근거 \\
    \midrule
    1 & \texttt{zeeman} 항은 선택 사항이다. 항이 없는 사이트는 장과 결합하지 않는다.
      $B\neq0$인데 항이 없는 사이트가 있으면 경고하고 결과에 기록한다.
      & 영행렬 항을 강제하면 모델 작성 부담만 늘고, 계산 비용 차이는 대각화에 비해 무시할 수 있다. \\
    2 & $H_Z=-\mu_B B_\alpha g_{\alpha\beta}S_\beta$: $g_{\alpha\beta}$는 스핀 성분 $\beta$를 장 성분 $\alpha$에 연결한다.
      & Einstein 합 규약에서 자연스럽고, 비대칭 $g$의 전치 모호성을 없앤다. \\
    3 & (D20으로 대체) 에너지 단위가 \texttt{relative}인 모델은 $\mu_B=1$로 두고 $B$를 에너지 단위 값으로 해석한다.
      & 상대 단위 모델에서 환산 근거를 추정하지 않는다. \\
    4 & 상태 검증(사이트 대응, 단위 벡터, 초격자 유효성, 계산계 정합, 주기 최소성)과
      기준 상태 진단(정상성, 안정성)을 나눈다. 방향은 단위 벡터로 저장한다.
      & 정상성·안정성은 모델과 $B$에 따라 달라진다. 접힌 밴드는 Chern 수를 모호하게 한다.
      각도 표현은 극점에서 $\phi$가 정의되지 않는다. \\
    \bottomrule
  \end{tabularx}
  \par\medskip
  \slidesource{2026-09-29 세부 규칙 결정}{contract}
\endgroup


\subsection{2026-09-29 당시의 상태 탐색: DE 뒤 세 경로}\label{state-search-current}
\begingroup% frame-local settings, as in Beamer
  \framesubtitle{기록용. 지금은 \texttt{classical}만 남았다(\texttt{quantum}은 D18·D22, \texttt{MAGSWT}는 D27에서 삭제).}
  \small
  \begin{tabularx}{\textwidth}{@{}p{1.7cm}YP{2.2cm}Y@{}}
    \toprule
    경로 & DE 뒤 단계 & 결과 method & 한계 \\
    \midrule
    \texttt{classical} & DE 각도에서 $E_{\mathrm{qm}}$만 계산 & \texttt{DE}
      & 고전 축퇴를 가르지 않음 \\
    \texttt{MAGSWT} & 모든 azimuth에 $k\pi/6$ ($k=1,\dots,6$)을 더한 점과 DE 점에서 $E_{\mathrm{cl}}+E_{\mathrm{qm}}$ 비교
      & \texttt{DE(MAGSWT)}, \texttt{MAGSWT}
      & $z$축 회전만 가정; \texttt{tl\_angle}의 $\theta$ 성분은 회전이 아님; 격자 간격이 분해능 \\
    \texttt{quantum} & L-BFGS-B로 모든 자유 각도의 $E_{\mathrm{cl}}+E_{\mathrm{qm}}$ 최소화, 무작위 교란 재시작 2회
      & \texttt{DE+BFGS}, \texttt{DE}
      & manifold 제약 없음; 교란은 전역 \texttt{np.random} \\
    \bottomrule
  \end{tabularx}
  \par\medskip
  \supportingtext{DE는 \texttt{seed=42}로 고정되어 결정론적이며, \texttt{angle\_setting}으로 일부 각도를 고정할 수 있다.
  \texttt{quantum}의 fallback·배열 변경 결함은 \texttt{d43d8d4}에서 고쳤다(부록 A). 영점 선택은 \texttt{select\_on\_manifold}, 궤도 위 에너지 지형은 \texttt{orbit\_energy\_landscape}.}
  \slidesource{현재 \texttt{methods/optimization.py}}{implementation,selection}
\endgroup


\subsection{D17: 영점 에너지 선택은 고전 manifold 위에서 한다}\label{state-selection-rule}
\begingroup% frame-local settings, as in Beamer
  \framesubtitle{축은 해밀토니안의 대칭이 아니라 고전 에너지 $E_{\mathrm{cl}}$에서 정한다.}
  \begin{columns}[T,onlytextwidth]
    \begin{column}{0.46\textwidth}
      \begin{block}{manifold 위에서 비교하는 이유}
        \small
        \begin{itemize}
          \item manifold 밖에서는 토크가 0이 아니어서 보손 선형항이 남는다.
          \item $E_{\mathrm{cl}}=O(S^2)$, $E_{\mathrm{qm}}=O(S^0)$의 합을 자유롭게 최소화하면
            $O(1/S)$만큼 벗어나 LSWT 차수를 넘는다.
          \item 정확한 대칭 방향에서는 $E_{\mathrm{qm}}$도 일정하다.
            선택은 $E_{\mathrm{cl}}$만 일정한 accidental 방향에서 일어난다.
        \end{itemize}
      \end{block}
    \end{column}
    \begin{column}{0.50\textwidth}
      \begin{block}{탐색 절차}
        \small
        \begin{enumerate}
          \item DE 해에서 국소 접평면 좌표로 $E_{\mathrm{cl}}$ Hessian 계산
          \item 회전 생성자 $\hat n\times\mathbf S_\alpha$를 영공간에 사영해 축 $\hat n$ 결정.
            영공간 없음 $\to$ DE 반환; 회전 span 밖 $\to$ 경고;
            여러 축 $\to$ 축 지정 요구
          \item 궤도 $R_{\hat n}(\phi)\mathbf S(0)$: $\phi$ 격자 뒤 1차원 bounded 최소화
          \item 궤도 전체의 $E_{\mathrm{cl}}$ 일정성 검사; 선택 에너지와 곡률 $C_\phi$ 기록
        \end{enumerate}
      \end{block}
    \end{column}
  \end{columns}
  \par\medskip
  \supportingtext{$\phi=0$이 후보에 들어가므로 결과는 항상 정의되고 난수가 필요 없다.
  $C_\phi$는 pseudo-Goldstone gap $\Delta_{\mathrm{PG}}^2=C_\phi/\chi_z$로 이어진다.
  $(\theta,\phi)$ 좌표는 극점 사이트에 가짜 영모드를 만들므로 쓰지 않는다.}
  \slidesource{2026-09-29 결정; NBCP 노트 6장의 궤도 정의를 일반 축으로 확장}{contract,selection}
\endgroup


\subsection{D18: $E_{\mathrm{cl}}+E_{\mathrm{qm}}$의 제약 없는 최소화를 없앤다}\label{quantum-removal}
\begingroup% frame-local settings, as in Beamer
  \framesubtitle{표준 스핀파 전개의 전제를 깨고, 그 에너지는 LSWT 차수를 넘는다.}
  \[
    H = S^2 e_{\mathrm{cl}}(\theta_0) + H_1 + H_2 + H_3 + H_4 + \cdots,
    \qquad H_1 \propto S^{3/2},\ H_2 \propto S,\ H_3 \propto S^{1/2},\ H_4 \propto S^0
  \]
  \small
  \begin{tabularx}{\textwidth}{@{}P{3.3cm}Y@{}}
    \toprule
    문제 & 내용 \\
    \midrule
    전제가 깨짐 & 전개의 전제는 $\theta_0$가 고전 정상점이라 $H_1=0$인 것. 최소점은 $\delta\theta\sim O(1/S)$ 벗어나 $H_1=O(S^{1/2})$가 남는데, $E_{\mathrm{qm}}$은 $H_1$을 버린 영점 에너지다 \\
    차수 불일치 & 에너지 이득 $S^2(1/S)^2=O(S^0)$는 버린 $H_3$·$H_4$ 보정과 같은 차수; 조화 근사 에너지는 변분 상한도 아니다 \\
    선택 오염 & NBCP의 선택 에너지($2.8\times10^{-6}$, $1.6\times10^{-12}$)보다 후보별 $O(S^0)$ 이완 차이가 클 수 있다 \\
    \bottomrule
  \end{tabularx}
  \par\medskip
  \supportingtext{정당한 부분은 manifold 위의 선택(D17)과, 기준점을 옮기지 않고 계산하는 모멘트 방향의 $1/S$ 보정(필요 시 별도 관측량)이다.
  \texttt{quantum}은 2b단계에서 D17 구현과 함께 삭제했다(D22).}
  \slidesource{2026-09-29 결정}{contract}
\endgroup


\subsection{D17 수치 확인: 축은 $E_{\mathrm{cl}}$ Hessian에서 정해진다}\label{state-selection-check}
\begingroup% frame-local settings, as in Beamer
  \framesubtitle{NBCP Y 상태, $B=0.2$ T, $N=6$, 궤도 24점. 모든 경우 영공간은 1차원이고 회전 span 안에 있다.}
  \small
  \begin{tabularx}{\textwidth}{@{}YP{2.6cm}P{2.3cm}P{4.2cm}@{}}
    \toprule
    경우 & 찾은 축 & 궤도 $E_{\mathrm{cl}}$ 변동 & 궤도 $E_{\mathrm{qm}}$ 변동 \\
    \midrule
    정확한 U(1) & $\hat z$, 오차 $9\times10^{-9}$ & $7\times10^{-18}$ & $3.8\times10^{-16}$: 선택 없음 \\
    $J_{\mathrm{PD}}=0.010$ & $\hat z$, 오차 $2\times10^{-8}$ & $1\times10^{-17}$ & $2.83\times10^{-6}$: $\pi/3$ 간격 최소 6개 \\
    $J_\Gamma=0.010$ & $\hat z$, 오차 $1\times10^{-8}$ & $1\times10^{-17}$ & $3.2\times10^{-12}$: 수렴한 선택(다음 쪽) \\
    $J_{\mathrm{PD}}$, 문제 전체를 $R$로 회전 & $R\hat z$, 오차 $8\times10^{-9}$ & $2\times10^{-17}$ & $2.83\times10^{-6}$: 회전 전과 같음 \\
    같은 회전에서 azimuth 이동 $R_z$ & --- & $1.5\times10^{-2}$ & 의미 없음 \\
    \bottomrule
  \end{tabularx}
  \par\medskip
  DE 해(기울기 $2$--$3\times10^{-7}$)에서 출발해도 축 오차는 $1.3\times10^{-5}$ 이하였다.
  다만 최소 고윳값 $3.6\times10^{-7}$(다음 $6.7\times10^{-3}$)이 고정 허용오차로는 영모드로 판정되지 않은 경우가 있었다.
  \par\medskip
  \supportingtext{영공간 판정 기준과 선택 판정 규칙의 자료와 검토안은 다음 두 쪽에 있다.}
  \slidesource{\texttt{examples/nbcp\_y\_orbit\_axis\_check.py}}{selection}
\endgroup


\subsection{D17 판정 자료 1: 영공간은 정밀화 뒤 간격비로 판정한다}\label{state-selection-null-criteria}
\begingroup% frame-local settings, as in Beamer
  \framesubtitle{판정 자료 · DE 세 seed $\times$ 축퇴 4경우(12건)와 대조 2경우(6건). 대조: 장을 기울임, $20h$ 편극 상태.}
  \small
  \begin{tabularx}{\textwidth}{@{}YP{3.0cm}P{3.0cm}@{}}
    \toprule
    판정 규칙 & 축퇴 12건 판정 & 대조 6건 오판 \\
    \midrule
    DE 해에 고정 허용오차(최대 고윳값의 $10^{-6}$배) & 8건 & 0건 \\
    DE 해에 간격비 $w_{\min}/w_{\mathrm{next}}<10^{-3}$ & 12건(최대 $5.3\times10^{-5}$) & 0건(최소 $2\times10^{-2}$) \\
    고전 L-BFGS-B 정밀화 뒤 고정 허용오차 & 12건(간격비 최대 $3\times10^{-7}$) & 0건 \\
    \bottomrule
  \end{tabularx}
  \par\medskip
  정밀화는 $E_{\mathrm{cl}}$을 $1.7\times10^{-11}$ 이하로 바꾸고 축 오차를 $1.4\times10^{-5}$에서 $1.1\times10^{-7}$로 줄였다.
  편극 상태에서 사라지는 $z$ 생성자의 상대 특이값은 DE 해에서 $3$--$8\times10^{-6}$, 정밀화 뒤 $2$--$4\times10^{-8}$이다.
  \par\medskip
  \textbf{D19}: 이 표의 간격비·rank 값은 고정 상수 모드의 기본값이 된다. 기본 모드는 물리 근거로 판정한다(다음 쪽들).
  \par\medskip
  \supportingtext{기울인 장 대조군(간격비 $2\times10^{-2}$)과 축퇴 DE 해($5\times10^{-5}$)의 여유는 약한 고전 pinning에서 줄어든다.
  그래서 판정만이 아니라 간격비를 기록한다.}
  \slidesource{\texttt{examples/nbcp\_y\_orbit\_criteria\_check.py}}{criteria}
\endgroup


\subsection{D17 판정 자료 2: 선택 위치는 작은 mesh에서도 정해진다}\label{state-selection-mesh}
\begingroup% frame-local settings, as in Beamer
  \framesubtitle{검토안 · 궤도 위 $E_{\mathrm{qm}}$의 6배 조화 성분 진폭을 독립 경로의 수렴값 $\lambda_6$($N=96$)과 비교했다.}
  \small
  \begin{tabularx}{\textwidth}{@{}YP{2.7cm}P{1.6cm}P{1.6cm}P{1.6cm}@{}}
    \toprule
    경우 & $\lambda_6$ (meV/spin) & $N=3$ & $N=6$ & $N=12$ \\
    \midrule
    $J_{\mathrm{PD}}=0.010$ & $1.44\times10^{-6}$ & $+9.2\%$ & $+1.1\%$ & $+0.2\%$ \\
    $J_{\mathrm{PD}}=0.005$ & $1.19\times10^{-7}$ & $+4.6\%$ & $+0.6\%$ & $+0.1\%$ \\
    $J_\Gamma=0.010$ & $1.64\times10^{-12}$ & $+3.5\%$ & $0.00\%$ & $0.00\%$ \\
    $J_\Gamma=0.005$ & $2.53\times10^{-14}$ & $+3.3\%$ & $-0.05\%$ & $-0.01\%$ \\
    \bottomrule
  \end{tabularx}
  \par\medskip
  모든 경우와 $N\ge3$에서 최소 위치는 $\phi\equiv0\pmod{\pi/3}$이었다(차이 $7.5\times10^{-5}$ rad 이하).
  앞 쪽의 $J_\Gamma$ 변동 $3.2\times10^{-12}$는 수렴한 선택이다($2\lambda_6$). 24점 격자는 위치 분해능만 제한한다.
  \par\medskip
  \textbf{검토안}: 궤도 위 $E_{\mathrm{qm}}$은 격자 최솟값 대신 조화 성분 fit으로 위치와 진폭을 정한다.
  진폭이 부동소수 바닥(상대 $10^{-16}$ 수준)에 여유를 둔 값보다 작으면 ``선택 없음''으로 기록한다.
  \par\medskip
  \supportingtext{비교 기준은 \texttt{examples/nbcp\_y\_angular\_matching.py}의 기록이며, 이 쪽의 값은 \texttt{EnergyFunction} 경로다.}
  \slidesource{\texttt{examples/nbcp\_y\_orbit\_criteria\_check.py}}{criteria}
\endgroup


\subsection{D19: 판정 기준은 물리적 질문에 묶는다}\label{state-selection-criteria}
\begingroup% frame-local settings, as in Beamer
  \framesubtitle{고정 임계값은 NBCP Y 상태 하나로 맞춘 값이라 모델마다 다를 수 있다. 두 모드 중 선택한다.}
  \small
  \begin{tabularx}{\textwidth}{@{}P{2.0cm}P{3.2cm}YP{2.9cm}@{}}
    \toprule
    판정 & 물리적 질문 & 물리 근거 모드(기본) & 고정 상수 모드 \\
    \midrule
    영공간 & 이 방향은 평평한가 & 정밀화 뒤 토크로 추정한 수치 정확도 이하면 평평; 분해되면 궤도 방향 $C_{\mathrm{cl}}$ 대 $C_{\mathrm{qm}}$로 고전 고정·양자 선택·경쟁 & 간격비 $<10^{-3}$ \\
    선택 없음 & $E_{\mathrm{qm}}$이 변하지 않아야 하는가 & 궤도 회전이 모든 항의 대칭인지($R^{\mathsf T}JR=J$, zeeman) 모델에서 판정; 아니면서 진폭이 fit 잔차 수준이면 ``분해 불가'' & 진폭 $<\max(10\times$잔차$,\,10^3\epsilon|E_{\mathrm{qm}}|)$ \\
    생성자 rank & 회전이 상태를 움직이는가 & 모든 스핀이 $\hat n$과 공선인지; 허용오차는 추정한 각도 정확도의 배수 & 상대 특이값 $10^{-4}$ \\
    \bottomrule
  \end{tabularx}
  \par\medskip
  \supportingtext{물리 근거는 모두 가벼운 계산이다(토크, Hessian, 이미 계산한 $E_{\mathrm{qm}}$, 행렬 연산, 외적).
  고정 상수의 기본값은 \texttt{definitions/defaults.py}에 두고 실행별 값은 방법 설정으로 넘긴다. 어느 모드든 판정에 쓴 값을 결과에 기록한다.}
  \slidesource{2026-09-29 결정; 세부 계수는 부록 A의 2b단계 쪽}{contract}
\endgroup


\subsection{D28: 비교 규칙을 1-loop 유효 퍼텐셜로 바꾼다}\label{d28-principle}
\begingroup% frame-local settings, as in Beamer
  \framesubtitle{2026-09-30 채택·구현(\texttt{0e9247f}, 부록 A). 이전 규칙(곡률 비, 대역 0.1--10)은 근거가 없고 조화 차수가 다르면 $m^2$만큼 치우친다.}
  \small
  \[
    \Gamma(\phi)=\min_{\text{다른 좌표}}E_{\mathrm{cl}}\big|_{\phi}+E_{\mathrm{zp}}(\phi)+O(S^0)
  \]
  \begin{tabularx}{\textwidth}{@{}P{2.6cm}Y@{}}
    \toprule
    항목 & 내용 \\
    \midrule
    의미 & 부드러운 좌표 $\phi$의 1-loop 유효 퍼텐셜; 답은 $\Gamma$의 최소(Shender 1982, Henley 1989; 유사 Goldstone 갭: Rau--McClarty--Moessner 2018) \\
    완화 경로 & 각 $\phi$에서 궤도 방향만 고정하고 나머지를 고전 최소화. 궤도 방향 토크만 남고, 스핀에 수직인 구속장은 $H_2$를 바꾸지 않으므로 선형항을 버린 계산이 구속된 1-loop 계산 \\
    D17·D18과의 관계 & 고전 궤도가 평평하면 $\Gamma=\text{상수}+E_{\mathrm{zp}}$로 D17과 같다. $\phi$ 하나만 움직여 D18이 막은 딱딱한 좌표의 $O(1/S)$ 이완이 없다 \\
    적용 조건 & 단열 분리: $\Gamma$의 곡률 $\ll$ 딱딱한 모드의 최소 강성. 비를 결과에 기록 \\
    \bottomrule
  \end{tabularx}
  \par\medskip
  \supportingtext{$d\Gamma/d\phi=0$: 부드러운 방향의 1-loop 선형항 상쇄. 한계 $O(S^0)$, $T=0$. 세부(확정): $\Gamma$ 최소 반환; \texttt{competition}·\texttt{not\_flat} 폐지, \texttt{selected}(이동량 보고)·\texttt{not\_soft} 추가; 단열 비 $>0.1$ 경고.}
  \slidesource{D28 검토 자료}{potential}
\endgroup


\subsection{D28 검토 자료: 강체 궤도는 고전 고정을 과대평가했다}\label{d28-results}
\begingroup% frame-local settings, as in Beamer
  \framesubtitle{NBCP Y(0.2 T, $J_{\mathrm{PD}}$)와 V(1.4 T, $J_\Gamma$), 면내 장 성분 비 = 기울기, 72점, $N=6$.}
  \footnotesize\renewcommand{\arraystretch}{1.08}
  \begin{tabularx}{\textwidth}{@{}P{1.9cm}P{2.3cm}P{2.3cm}P{2.1cm}Y@{}}
    \toprule
    경우 & 고전 변동: 강체 / 완화 & 양자 변동(완화) & $\Gamma$ 최소$-$양자 최소 & 단열 비 \\
    \midrule
    Y, 0° & 0 / 0 & $2.9\times10^{-6}$ & $3.7\times10^{-7}$ (D17) & $1.1\times10^{-2}$ \\
    Y, 0.57° & -- / $5.9\times10^{-8}$ & $3.2\times10^{-6}$ & $-9.7\times10^{-4}$ & $1.1\times10^{-2}$ \\
    Y, 1.7° & $7.5\times10^{-6}$ / $5.3\times10^{-7}$ & $3.8\times10^{-6}$ & $-8.9\times10^{-3}$ & $1.1\times10^{-2}$ \\
    Y, 5.7° & -- / $5.8\times10^{-6}$ & $9.0\times10^{-6}$ & $+4.5\times10^{-2}$ & $2.8\times10^{-2}$ \\
    Y, 16.7° & -- / $5.1\times10^{-5}$ & $5.7\times10^{-5}$ & $+6.1\times10^{-2}$ & $1.7\times10^{-1}$ \\
    V, 0.57° & -- / $1.8\times10^{-4}$ & $5.8\times10^{-5}$ & 고전 최소 쪽($-3.3\times10^{-2}$) & $9.3\times10^{-3}$ \\
    V, 5.7° & 완화 실패 & -- & -- & $2.4\times10^{3}$ \\
    \bottomrule
  \end{tabularx}
  \par\smallskip
  \supportingtext{정정: 강체 궤도로 얻은 ``1.7°면 고전 고정이 이긴다''는 틀렸다(고전 변동을 14배 과대평가). Y는 면내 알짜 모멘트가 없어 고전 변동 $\propto$ 기울기$^2$, V는 선형이라 0.57°에서 이미 고전 우세. 에너지 단위 meV/spin, 이동 단위 rad.
  \textbf{구현 후 정정}: 이 표는 MAGSWT 정규화 제공자로 계산해 기울인 장의 $\Gamma$ 최소·단열 비가 오염되었다(Y 1.7°는 $-8.3\times10^{-3}$, Y 16.7° 단열 비는 0.023). 정정 값은 부록 A의 D28 구현 쪽.}
  \slidesource{\texttt{examples/nbcp\_effective\_potential\_check.py}}{potential}
\endgroup


\subsection{설계 문서와 부록을 갱신하는 방법}\label{record-management}
\begingroup% frame-local settings, as in Beamer
  \framesubtitle{변경 내용, 근거, 판정, 다음 작업을 함께 남긴다.}
  \begin{columns}[T,onlytextwidth]
    \begin{column}{0.48\textwidth}
      \begin{block}{내용의 소유 위치}
        \textbf{전달 규약 상세}\\
        Working-Pad의 검토 문서\\[0.7em]
        \textbf{개발 목표와 시스템 설계}\\
        \texttt{main.tex}와 \texttt{sections/}\\[0.7em]
        \textbf{수치 증거와 소스 식별}\\
        검증 JSON · \texttt{source-manifest.json}
      \end{block}
    \end{column}
    \begin{column}{0.48\textwidth}
      \begin{block}{문서의 갱신 순서}
        \begin{enumerate}
          \item 본문의 목표·구조·규약을 갱신한다.
          \item 구현·시험 기록을 부록에 연결한다.
          \item 계획·결정·열린 이슈를 갱신한다.
          \item 기록일·소스 manifest를 갱신한다.
          \item PDF를 빌드하고 전 페이지를 확인한다.
        \end{enumerate}
      \end{block}
    \end{column}
  \end{columns}
  \medskip
  \supportingtext{화면에 보이는 PDF는 편집 원본에서 생성한다.
  과거 검증 수치를 새 코드의 결과로 자동 승계하지 않는다.}
  \slidesource{설계 문서의 관리 방식; 템플릿은 Knowledge Factory에서 재사용}{contract,theme}
\endgroup


\subsection{근거 문서와 템플릿}\label{sources}
\begingroup% frame-local settings, as in Beamer
  \framesubtitle{긴 경로와 정확한 해시는 source-manifest.json에서 관리한다.}
  \small
  \begin{thebibliography}{12}
    \bibitem{contract}
      모델 간 공통 SpinModel 전달 규약, 2026-09-23 작성 · 2026-09-29 갱신.
      \newblock \texttt{archive/government/Working-Pad/idea-proposals/} · in-review.
    \bibitem{plan}
      2D Spin-System Toolkit 1차 개발 계획, 2026-09-29.
      \newblock \texttt{archive/government/Working-Pad/idea-proposals/} · in-review · 범위·패키지·벤치마크·단계.
    \bibitem{extraction}
      NBCP 모델 추출 검증 기록, 2026-09-23.
      \newblock \texttt{code-space/model/nbcp/verification/} · 모델·해밀토니안·탐색·pytest 기록.
    \bibitem{implementation}
      현재 \texttt{SpinSystem}, 격자, 고전 에너지, LSWT 해밀토니안 구현.
      \newblock \texttt{code-space/spintoolkit/} · 작업 공간 소스 대조.
    \bibitem{geometry}
      NBCP Y SOC 조건 검사: \texttt{geometry\_audit}.
      \newblock \texttt{examples/nbcp\_y\_soc\_conditions.py} · 코드 관찰 근거.
    \bibitem{theme}
      Knowledge Factory Beamer, 로컬 템플릿 스냅샷.
      \newblock 선언 버전 v0.2.3 · note mode · 복사한 theme 내용은 수정하지 않음.
    \bibitem{layout}
      기능별 패키지 이행 검증 기록, 2026-09-23.
      \newblock \texttt{archive/2026-10-08-development-beamer/verification/} · 본문·수치·호환성·pytest 대조.
    \bibitem{rename}
      패키지 이름 변경 검증 기록, 2026-09-29.
      \newblock \texttt{archive/2026-10-08-development-beamer/verification/} · 테스트·수치·호환성·설치 대조.
    \bibitem{stage1}
      1단계 공통 자료형 검증 기록, 2026-09-29.
      \newblock \texttt{archive/2026-10-08-development-beamer/verification/} · 정확값·교차 확인·규약 사례·회귀 대조.
    \bibitem{selection}
      상태 선택 검증 기록, 2026-09-29.
      \newblock \texttt{archive/2026-10-08-development-beamer/verification/} · \texttt{quantum} 경로 수정과 D17 궤도 축 수치 확인.
  \end{thebibliography}
\endgroup


\subsection{근거 문서 (계속)}\label{sources-continued}
\begingroup% frame-local settings, as in Beamer
  \framesubtitle{긴 경로와 정확한 해시는 source-manifest.json에서 관리한다.}
  \small
  \begin{thebibliography}{12}
    \setcounter{enumiv}{10}
    \bibitem{criteria}
      D17 판정 기준 자료, 2026-09-29.
      \newblock \texttt{archive/2026-10-08-development-beamer/verification/} · 영공간 판정 규칙 비교와 선택의 mesh 의존성.
    \bibitem{stage2}
      2단계 NBCP 연결 검증 기록, 2026-09-29.
      \newblock \texttt{archive/2026-10-08-development-beamer/verification/} · $\mathsf H(\mathbf k)$ 원소 단위 대조, 왕복 변환, 에너지.
    \bibitem{stage4a}
      4a LSWT 진입점·바닥상태 검증 기록, 2026-09-30.
      \newblock \texttt{archive/2026-10-08-development-beamer/verification/} · 기존 솔버 재현, 해석값, QMC·DMRG.
    \bibitem{stage4b}
      4b 영모드·유한 온도 검증 기록, 2026-09-30.
      \newblock \texttt{archive/2026-10-08-development-beamer/verification/} · 영모드 판정, ED 저온 대조.
    \bibitem{stage4c}
      4c 구조인자 검증 기록, 2026-09-30.
      \newblock \texttt{archive/2026-10-08-development-beamer/verification/} · ED 무게, 합 규칙, 기존 코드 진단.
    \bibitem{cleanup}
      기존 코드 정리 기록, 2026-09-30.
      \newblock \texttt{archive/2026-10-08-development-beamer/verification/} · MAGSWT 격자 탐색·\texttt{correlations.py} 삭제와 대체.
    \bibitem{potential}
      D28 검토 자료: 1-loop 유효 퍼텐셜, 2026-09-30.
      \newblock \texttt{archive/2026-10-08-development-beamer/verification/} · 완화 경로, D17 극한, 기울인 장.
    \bibitem{d28impl}
      D28 구현 검증 기록, 2026-09-30.
      \newblock \texttt{archive/2026-10-08-development-beamer/verification/} · 시제품 대조, 영점 제공자, mesh 수렴, 2b 재실행; \texttt{d28-provider-projection-2026-09-30.json} · $k=0$ 구속 모드 투영.
    \bibitem{stage6b}
      6b 검증 기록, 2026-09-30.
      \newblock \texttt{archive/2026-10-08-development-beamer/verification/stage6b-luttinger-tisza-2026-09-30.json} · \texttt{examples/luttinger\_tisza\_check.py}.
    \bibitem{stage6c}
      6c 검증 기록, 2026-09-30.
      \newblock \texttt{archive/2026-10-08-development-beamer/verification/stage6c-classical-search-2026-09-30.json} · \texttt{examples/classical\_search\_check.py}.
    \bibitem{stage4d}
      4d 검증 기록, 2026-09-30.
      \newblock \texttt{archive/2026-10-08-development-beamer/verification/stage4d-band-structure-2026-09-30.json} · \texttt{examples/lswt\_stage4d\_check.py}.
    \bibitem{stage7}
      7단계 검증 기록, 2026-10-01.
      \newblock \texttt{archive/2026-10-08-development-beamer/verification/stage7-visualization-2026-10-01.json} · \texttt{examples/visualization\_stage7\_check.py}.
    \bibitem{stage2b}
      2b단계 상태 선택 검증 기록, 2026-09-29.
      \newblock \texttt{archive/2026-10-08-development-beamer/verification/} · NBCP Y·V 선택, 대조군, 벤치마크, MAGSWT 재현.
    \bibitem{stage3}
      3단계 토러스 전개·최소 ED 검증 기록, 2026-09-30.
      \newblock \texttt{archive/2026-10-08-development-beamer/verification/} · 1-마그논 대 LSWT, 독립 구현 대조.
  \end{thebibliography}
\endgroup


\subsection{근거 문서 (5단계 이후)}\label{sources-stage5}
\begingroup% frame-local settings, as in Beamer
  \framesubtitle{긴 경로와 정확한 해시는 source-manifest.json에서 관리한다.}
  \small
  \begin{thebibliography}{24}
    \setcounter{enumiv}{20}
    \bibitem{stage5a}
      5a 검증 기록, 2026-09-30.
      \newblock \texttt{archive/2026-10-08-development-beamer/verification/stage5a-berry-chern-2026-09-30.json} · \texttt{examples/lswt\_stage5a\_check.py}.
    \bibitem{stage5b}
      5b 검증 기록, 2026-09-30.
      \newblock \texttt{archive/2026-10-08-development-beamer/verification/stage5b-thermal-hall-2026-09-30.json} · \texttt{examples/lswt\_stage5b\_check.py}.
    \bibitem{stage5c}
      5c 검증 기록, 2026-09-30.
      \newblock \texttt{archive/2026-10-08-development-beamer/verification/stage5c-hall-kernel-2026-09-30.json}.
    \bibitem{stage5d}
      5d 검증 기록, 2026-09-30.
      \newblock \texttt{archive/2026-10-08-development-beamer/verification/stage5d-adaptive-integration-2026-09-30.json} · \texttt{examples/lswt\_stage5d\_check.py}.
    \bibitem{stage6b-quarter}
      6b 보정 검증 기록, 2026-09-30.
      \newblock \texttt{archive/2026-10-08-development-beamer/verification/stage6b-quarter-wave-vector-2026-09-30.json} · \texttt{test\_quarter\_wave\_vector\_accepts\_the\_up\_up\_down\_down\_state}.
  \end{thebibliography}
\endgroup

